\documentclass[11pt,reqno]{amsart}
\usepackage[utf8]{inputenc}
\usepackage[T1]{fontenc}

\usepackage[margin=1.15in]{geometry}
\usepackage{amsmath,amssymb,amsthm,mathtools}
\usepackage{tikz-cd}
\usepackage{enumitem}
\usepackage{booktabs}
\usepackage[colorlinks=true,linkcolor=blue!50!black,citecolor=green!40!black,urlcolor=blue!50!black]{hyperref}
\usepackage[capitalise,nameinlink,noabbrev]{cleveref}

\setlist[enumerate]{label=\textup{(\arabic*)}}
\setlength{\emergencystretch}{3em}

% ---------------------------------------------------------------
% Theorem environments
% ---------------------------------------------------------------
\theoremstyle{plain}
\newtheorem{theorem}{Theorem}[section]
\newtheorem{proposition}[theorem]{Proposition}
\newtheorem{lemma}[theorem]{Lemma}
\newtheorem{corollary}[theorem]{Corollary}
\newtheorem{conjecture}[theorem]{Conjecture}
\newtheorem{question}[theorem]{Question}
\newtheorem{maintheorem}{Theorem}
\renewcommand{\themaintheorem}{\Alph{maintheorem}}
\theoremstyle{definition}
\newtheorem{definition}[theorem]{Definition}
\newtheorem{example}[theorem]{Example}
\newtheorem{setting}[theorem]{Setting}
\newtheorem{notation}[theorem]{Notation}
\theoremstyle{remark}
\newtheorem{remark}[theorem]{Remark}
\newtheorem*{intuition}{Intuition}

\crefname{maintheorem}{Theorem}{Theorems}
\crefname{setting}{Setting}{Settings}
\crefname{notation}{Notation}{Notations}
\crefname{question}{Question}{Questions}
\crefname{conjecture}{Conjecture}{Conjectures}

% ---------------------------------------------------------------
% Macros
% ---------------------------------------------------------------
\newcommand{\kk}{\Bbbk}
\DeclareMathOperator{\Hom}{Hom}
\DeclareMathOperator{\End}{End}
\DeclareMathOperator{\Ext}{Ext}
\DeclareMathOperator{\Filt}{Filt}
\DeclareMathOperator{\add}{add}
\DeclareMathOperator{\real}{real}
\DeclareMathOperator{\gldim}{gl.dim}
\DeclareMathOperator{\pd}{pd}
\DeclareMathOperator{\soc}{soc}
\DeclareMathOperator{\rad}{rad}
\DeclareMathOperator{\tr}{tr}
\DeclareMathOperator{\im}{Im}
\DeclareMathOperator{\cone}{cone}
\DeclareMathOperator{\RHom}{\mathbf{R}Hom}
\newcommand{\modd}{\operatorname{mod}}
\newcommand{\proj}{\operatorname{proj}}
\newcommand{\Db}{\mathcal{D}^b}
\newcommand{\Kb}{K^b}
\newcommand{\thick}[1]{\langle #1\rangle}
\newcommand{\Pn}[2]{\mathbb{P}^{#1}[#2]}
\newcommand{\nablabar}{\overline{\nabla}}
\newcommand{\Deltabar}{\overline{\Delta}}
\newcommand{\cD}{\mathcal{D}}
\newcommand{\cA}{\mathcal{A}}
\newcommand{\cH}{\mathcal{H}}
\newcommand{\cT}{\mathcal{T}}
\newcommand{\cC}{\mathcal{C}}
\newcommand{\cF}{\mathcal{F}}
\newcommand{\cN}{\mathcal{N}}
\newcommand{\cM}{\mathcal{M}}
\newcommand{\cQ}{\mathcal{Q}}
\newcommand{\cU}{\mathcal{U}}
\newcommand{\cV}{\mathcal{V}}
\newcommand{\cO}{\mathcal{O}}
\newcommand{\pH}{{}^{p}\!H}

\title[Local-exceptional sequences and stratified algebras]{A stratified Bodzenta--Bondal criterion: local-exceptional sequences, glued t-structures and stratified algebras}
\date{October 2026}

\begin{document}

\begin{abstract}
Bodzenta and Bondal characterised highest weight categories as hearts of t-structures glued along full exceptional sequences, under a vanishing condition on negative extensions within the sequence and within its left dual. We prove a stratified analogue of this criterion.
\begin{itemize}
\item Exceptional objects are replaced by \emph{local-exceptional} objects. These have no self-extensions in non-zero degrees, but their endomorphism algebra may be an arbitrary local algebra.
\item The left dual sequence is replaced by a \emph{Hom-dual} sequence of ``simple-like'' objects, which generates the ambient category.
\end{itemize}
Assume that the dual sequence has a right dual in the sense of mutations; this is automatic for exceptional sequences. Then the glued heart is the module category of a standardly stratified algebra, with the two sequences as standard and proper costandard objects, if and only if neither sequence has negative extensions. In that case the ambient category is the bounded derived category of the heart.

The proof rests on a gluing theorem for admissible filtrations whose strata carry bounded t-structures with local hearts. Conversely, the standard modules of a standardly stratified algebra form a local-exceptional sequence. Their Hom-dual partners, the proper costandard modules, lie outside the subcategory generated by the standard modules unless the algebra is quasi-hereditary. This explains why, unlike in the exceptional case, the dual sequence cannot be obtained by mutation.

As an application we discuss $\Pn nk$-like objects. For $k=0$ they are local-exceptional. For $k\ge1$ they generate strata with local hearts, and gluing along sequences of them yields highest weight categories for $k\ge3$ and standardly stratified categories for $k=1$, $n\ge3$. When $n\geq1$, the ambient category is never the derived category of the heart.
\end{abstract}

\maketitle

\setcounter{tocdepth}{1}
\tableofcontents

% ===============================================================
\section{Introduction}
% ===============================================================

\subsection{Background}
Highest weight categories, introduced by Cline, Parshall and Scott \cite{CPS88}, are precisely the module categories of quasi-hereditary algebras. Their bounded derived categories carry two distinguished full exceptional sequences: the standard objects $(\Delta(\lambda_1),\dots,\Delta(\lambda_n))$, and the costandard objects, which form the left dual sequence $(\nabla(\lambda_n),\dots,\nabla(\lambda_1))$.

Bodzenta and Bondal \cite{BB26} proved that this characterises highest weight categories, in two ways:
\begin{itemize}
\item The heart of the t-structure glued along a full exceptional sequence is highest weight, with the given sequence and its left dual as standard and costandard objects, if and only if neither sequence has negative extensions (\cref{thm:BB-3.1}).
\item Highest weight structures on an abelian category $\cA$ correspond bijectively to pairs of dual full exceptional sequences in $\Db(\cA)$ consisting of objects of $\cA$ (\cref{thm:BB-3.8}).
\end{itemize}

\begin{intuition}
A full exceptional sequence is the same as an admissible filtration of $\cD$ whose subquotients (the \emph{strata}) are all equivalent to $\Db(\kk)$. Each stratum has, up to shift, a unique bounded t-structure. Its heart is $\modd\kk$, whose unique simple object $\kk$ is at the same time projective. The standard object $E_s$ is the extension by zero $j_{s!}(\kk)$ of this object, and the costandard object $F_s$ is its extension $j_{s*}(\kk)$.
\end{intuition}

\subsection{The question}
\emph{Standardly stratified} and \emph{properly stratified} algebras \cite{CPS96,Dla96,AHLU00} generalise quasi-hereditary algebras. The endomorphism algebras $B_i=\End(\Delta(i))$ of the standard modules are allowed to be arbitrary local algebras instead of $\kk$. We address the following question.

\begin{question}\label{q:main}
Is there a generalisation of dual pairs of exceptional sequences which recovers standardly stratified algebras, in the same way as dual pairs of exceptional sequences recover quasi-hereditary algebras?
\end{question}

The question arose from a more specific one, concerning $\Pn nk$-like objects \cite{HT06,Kru18,HK19}. These are objects $P$ with $\End^\bullet(P)\cong\kk[t]/(t^{n+1})$ and $\deg t=k$, which generalise exceptional objects (the case $n=0$) in a different direction. It turns out that the generalisation relevant for stratified algebras is a different one. $\Pn nk$-like objects are treated as an application in \cref{sec:P-objects}.

The natural candidates for standard objects are \emph{local-exceptional} objects (\cref{def:local-exc}): objects $E$ with $\Hom(E,E[l])=0$ for $l\ne0$ and $\End(E)$ local. Standard modules of standardly stratified algebras are of this kind (\cref{mthm:C}). Two observations show that the passage from exceptional to local-exceptional sequences is not formal.
\begin{itemize}
\item \emph{There is nothing to glue inside the category generated by the sequence.} The category generated by a local-exceptional object $E$ is $\Kb(\proj B)$, where $B=\End(E)$. It contains no object playing the role of the simple $B$-module (\cref{prop:local-exc}), and in important cases it has no non-trivial t-structure at all \cite{HJY13,Bonf26}.
\item \emph{The left dual is the wrong partner.} The left dual of a local-exceptional sequence pairs with it via $\Hom^\bullet(E_s,F_s)\cong B_s$ rather than $\kk$. The correct partners satisfy $\Hom^\bullet(E_s,N_s)\cong\kk$; for stratified algebras they are the proper costandard modules. These lie outside the subcategory generated by the $E_s$ unless $B_s=\kk$ (\cref{mthm:C}).
\end{itemize}
Consequently the dual sequence must be part of the data, and the ambient category must be generated by it.

\subsection{Main results}

\begin{maintheorem}[\cref{thm:B}]\label{mthm:B}
Let $\cD$ be a Hom-finite idempotent complete triangulated category, let $(E_1,\dots,E_m)$ be a local-exceptional sequence, and let $(N_m,\dots,N_1)$ be a sequence of objects such that:
\begin{itemize}
\item[\textup{(HD1)}] $\Hom^\bullet(E_s,N_t)\cong\delta_{st}\kk$, concentrated in degree $0$;
\item[\textup{(HD2)}] $\Hom^\bullet(N_s,N_t)=0$ for $s<t$; moreover $\End(N_s)=\kk$ and $\Hom(N_s,N_s[l])=0$ for $l<0$;
\item[\textup{(HD3)}] $\cD=\thick{N_1,\dots,N_m}$;
\item[\textup{(HD4)}] for each $s$ there are an object $\Deltabar_s$ and a morphism $\Deltabar_s\to N_s$ whose cone lies in $\thick{N_1,\dots,N_{s-1}}$, such that $\Hom^\bullet(\Deltabar_s,N_t)=0$ for all $t<s$.
\end{itemize}
Then $\cD_s=\thick{N_1,\dots,N_s}$ is an admissible filtration whose strata carry natural bounded t-structures with local hearts. Let $\cA$ be the heart of the glued t-structure. Then $\cA$ is standardly stratified, with standard objects $E_s$ and proper costandard objects $N_s$, if and only if
\[\Hom(E_s,E_t[l])=0=\Hom(N_s,N_t[l])\qquad\text{for all $s,t$ and all $l<0$.}\]
In this case:
\begin{itemize}
\item $\End(\Delta(s))\cong\End(E_s)$;
\item $\cA$ is quasi-hereditary if and only if all $E_s$ are exceptional;
\item if $\cD$ is algebraic, then $\real\colon\Db(\cA)\to\cD$ is an equivalence.
\end{itemize}
\end{maintheorem}

Condition (HD4) says that $(N_m,\dots,N_1)$ has a \emph{right dual} $(\Deltabar_1,\dots,\Deltabar_m)$ in the sense of mutations. Equivalently, each $\thick{N_1,\dots,N_{s-1}}$ is left admissible in $\thick{N_1,\dots,N_s}$ (\cref{rem:HD4}). The objects $\Deltabar_s$ play the role of proper standard objects. For an exceptional sequence, the left dual is a Hom-dual sequence whose right dual is the original sequence (\cref{prop:exc-dual}), and \cref{mthm:B} reduces to \cref{thm:BB-3.1}.

\Cref{cor:BB38} gives the corresponding analogue of \cref{thm:BB-3.8}. It is a bijection between \emph{dualisable} standardly stratified structures, i.e.\ those whose standard and proper costandard objects extend to a Hom-dual triple, and pairs of sequences of objects of $\cA$ that extend to Hom-dual triples.
\begin{itemize}
\item Quasi-hereditary structures are dualisable.
\item A standardly stratified structure is dualisable, with the proper standard modules as right dual, as soon as injective modules have $\nablabar$-filtrations and $\Ext^\bullet(\Deltabar(s),\nablabar(t))=0$ for $t<s$ (\cref{prop:dualisable}).
\end{itemize}
Whether every standardly stratified structure is dualisable is open (\cref{q:HD4}).

\Cref{mthm:B} is deduced from a gluing theorem for filtrations whose strata are \emph{local}: each stratum carries a bounded t-structure whose heart is a length category with a unique simple object (\cref{def:local-stratum}).

\begin{maintheorem}[\cref{thm:A}]\label{mthm:A}
Let $\cD$ be a Hom-finite triangulated category with an admissible filtration $0=\cD_0\subset\cD_1\subset\dots\subset\cD_m=\cD$ whose strata $\cT_s=\cD_s/\cD_{s-1}$ are local. Suppose that each simple object $S_s$ has a projective cover $\Pi_s$ in the heart of $\cT_s$. Put $\Delta_s=j_{s!}\Pi_s$ and $\nablabar_s=j_{s*}S_s$, and let $\cA$ be the glued heart. Assume that $\Delta_s,\nablabar_s\in\cA$ for all $s$ (equivalently, neither family has negative extensions), and that $\Hom_{\cT_s}(\Pi_s,S_s[2])=0$ for all $s$. Then:
\begin{enumerate}
\item $\cA\simeq\modd\Lambda$ for a standardly stratified algebra $\Lambda$, with standard modules $\Delta_s$, proper costandard modules $\nablabar_s$, and $\End(\Delta_s)\cong\End(\Pi_s)$.
\item If $\cD$ is algebraic, then $\real\colon\Db(\cA)\to\cD$ is an equivalence if and only if $\Hom_{\cT_s}(\Pi_s,S_s[l])=0$ for all $s$ and all $l\ge1$.
\end{enumerate}
\end{maintheorem}

\begin{maintheorem}[\cref{thm:C}]\label{mthm:C}
Let $\Lambda$ be a standardly stratified algebra, and put $B_i=\End(\Delta(i))$.
\begin{enumerate}
\item The standard modules form a full local-exceptional sequence in $\Kb(\proj\Lambda)$ with endomorphism algebras $B_i$.
\item $\Hom^\bullet(\Delta(i),\nablabar(i))\cong\kk$, and $\Hom^\bullet(\Delta(j),\nablabar(i))=0$ for $j>i$.
\item If $B_i\neq\kk$, then $\nablabar(i)\notin\Kb(\proj\Lambda)$.
\item $\gldim\Lambda<\infty$ if and only if $\Lambda$ is quasi-hereditary.
\end{enumerate}
\end{maintheorem}

\subsection{Method: induction on strata}
As in the theory of quasi-hereditary algebras, all proofs are by induction, peeling off one stratum at a time.
\begin{itemize}
\item \emph{Abelian side.} \Cref{lem:strat-induction} characterises standardly stratified categories inductively. The trace of the projective cover of the maximal simple object in every projective object must be projective (a ``standardly stratifying'' ideal), and the quotient category must again be standardly stratified. For quasi-hereditary algebras this is the characterisation by heredity chains (\cref{rem:heredity}).
\item \emph{Triangulated side.} \Cref{prop:step} glues one more local stratum onto a category whose heart is already known to be standardly stratified, and verifies the trace condition. \Cref{lem:peel} shows that the hypotheses of \cref{mthm:B} are inherited by $\thick{N_1,\dots,N_{m-1}}$, and that the last stratum is local of projective type.
\end{itemize}
\Cref{mthm:A,mthm:B} then follow by induction on the number of strata.

\subsection{Application: \texorpdfstring{$\mathbb{P}^n[k]$}{Pn[k]}-like objects}
In \cref{sec:P-objects} we specialise to $\Pn nk$-like objects, using Bonfert's formality and classification results \cite{Bonf26}.
\begin{itemize}
\item \emph{$k=0$.} $\Pn n0$-like objects are exactly the local-exceptional objects with endomorphism algebra $\kk[t]/(t^{n+1})$, so \cref{mthm:B,mthm:C} apply directly.
\item \emph{$k\ge1$.} A $\Pn nk$-like object generates a local stratum, whose heart we compute. Gluing along semi-orthogonal sequences of such objects (\cref{thm:Pseq}) yields:
\begin{itemize}
\item highest weight categories when all $k_s\ge3$;
\item standardly stratified categories with dual-number endomorphism algebras when $k_s=1$ and $n_s\ge3$;
\item categories that are not module categories when $k_s=n_s=1$.
\end{itemize}
In all cases with $n_s\ge1$, the ambient category is not the derived category of the heart. The case $k=2$, which contains the Huybrechts--Thomas $\mathbb{P}^n$-objects, is not decided (\cref{q:k2}). \Cref{tab:summary} summarises the outcome.
\end{itemize}

\subsection*{Sources and status of the statements}
Every statement attributed to the literature is quoted with its statement, in our notation. All statements without a citation are proved here, except those labelled \emph{Question}. From Bonfert's preprint \cite{Bonf26} we use the formality and classification results quoted in \cref{sec:P-objects}. Our computation of the hearts of the categories generated by $\Pn nk$-like objects (\cref{prop:heart-P}) overlaps with \cite[Lemmas~3.11 and~4.15]{Bonf26}, and the argument of \cref{lem:recognise-proj} is that of \cite[Lemma~3.12]{Bonf26}. We include proofs in the form needed here.

\subsection*{Organisation}
\begin{itemize}
\item \Cref{sec:prelim} collects definitions and the results we use.
\item \Cref{sec:local} introduces local-exceptional objects and local strata.
\item \Cref{sec:gluing} proves the gluing theorem (\cref{mthm:A}).
\item \Cref{sec:criterion} proves the stratified Bodzenta--Bondal criterion (\cref{mthm:B}).
\item \Cref{sec:stratified} treats standardly stratified algebras (\cref{mthm:C}) and the converse.
\item \Cref{sec:P-objects} contains the application to $\Pn nk$-like objects.
\item \Cref{sec:questions} lists open questions.
\end{itemize}

% ===============================================================
\section{Preliminaries}\label{sec:prelim}
% ===============================================================

\subsection{Conventions}
Throughout, $\kk$ is an algebraically closed field. Triangulated categories are $\kk$-linear, \emph{Hom-finite} (all $\Hom(X,Y)$ are finite-dimensional) and idempotent complete; in particular they are Krull--Schmidt. We use the following notation.
\begin{itemize}
\item For objects $X,Y$ we write $\Hom^l(X,Y):=\Hom(X,Y[l])$ and $\Hom^\bullet(X,Y):=\bigoplus_{l\in\mathbb{Z}}\Hom^l(X,Y)$. Thus $\End^\bullet(X)$ is a graded algebra under composition.
\item For a class $\mathcal S$ of objects, $\thick{\mathcal S}$ is the smallest thick subcategory (triangulated and closed under direct summands) containing $\mathcal S$.
\item $\Filt(\mathcal S)$ is the smallest full subcategory containing $0$ and $\mathcal S$ and closed under extensions: if $X\to Y\to Z\to X[1]$ is a triangle with $X,Z\in\Filt(\mathcal S)$, then $Y\in\Filt(\mathcal S)$.
\item We put $\mathcal S^\perp=\{Y:\Hom^\bullet(S,Y)=0\ \forall S\in\mathcal S\}$, and define ${}^\perp\mathcal S$ dually.
\end{itemize}
A triangulated category is \emph{algebraic} if it is equivalent to the homotopy category $H^0$ of a pretriangulated dg category (equivalently, to the stable category of a Frobenius exact category).

Modules are finite-dimensional right modules, and $\modd\Lambda$ denotes their category; $\Kb(\proj\Lambda)\subset\Db(\modd\Lambda)$ is the subcategory of perfect complexes. In an abelian length category, $[X:L]$ denotes the composition multiplicity of a simple object $L$ in $X$. For a set $\Theta$ of objects, $\cF(\Theta)$ is the full subcategory of objects admitting a finite filtration with subquotients isomorphic to objects in $\Theta$.

\subsection{t-structures and hearts}

\begin{definition}[{\cite[D\'ef.~1.3.1]{BBD82}}]
A \emph{t-structure} on $\cD$ is a pair $(\cD^{\le0},\cD^{\ge0})$ of strictly full subcategories such that, with $\cD^{\le n}=\cD^{\le0}[-n]$ and $\cD^{\ge n}=\cD^{\ge0}[-n]$:
\begin{enumerate}
\item $\cD^{\le0}[1]\subset\cD^{\le0}$ and $\cD^{\ge0}[-1]\subset\cD^{\ge0}$;
\item $\Hom(X,Y)=0$ for $X\in\cD^{\le0}$ and $Y\in\cD^{\ge1}$;
\item every $X$ fits into a triangle $X'\to X\to X''\to X'[1]$ with $X'\in\cD^{\le0}$ and $X''\in\cD^{\ge1}$.
\end{enumerate}
The \emph{heart} $\cA=\cD^{\le0}\cap\cD^{\ge0}$ is an abelian category, with cohomology functors $H^l\colon\cD\to\cA$. The t-structure is \emph{bounded} if every object lies in $\cD^{\le a}\cap\cD^{\ge b}$ for some $a,b$. It has \emph{length heart} if $\cA$ is a finite length category with finitely many isomorphism classes of simple objects.
\end{definition}

Short exact sequences in $\cA$ are exactly the triangles in $\cD$ all of whose vertices lie in $\cA$. Consequently $\cA$ is closed under extensions in $\cD$.

\begin{proposition}[{\cite[Rem.~3.1.17]{BBD82}}]\label{prop:ext-comparison}
Let $\cA$ be the heart of a t-structure on $\cD$, and let $X,Y\in\cA$. The natural maps $\Ext^l_\cA(X,Y)\to\Hom_\cD(X,Y[l])$, which send a Yoneda extension to the composite of the corresponding connecting morphisms, are bijective for $l\le1$ and injective for $l=2$. They are compatible with Yoneda product and composition.
\end{proposition}

\begin{theorem}[{Realisation; \cite[\S3.1]{BBD82}, \cite{Bei87}, \cite{Let24}}]\label{thm:realisation}
Let $\cD$ be algebraic (more generally, the base of a filtered triangulated category), and let $\cA$ be the heart of a bounded t-structure on $\cD$.
\begin{enumerate}
\item There is an exact functor $\real\colon\Db(\cA)\to\cD$ which restricts to the inclusion on $\cA$ and induces the maps of \cref{prop:ext-comparison} on extension groups.
\item $\real$ is fully faithful if and only if the following \emph{effaceability} condition holds: for all $X,Y\in\cA$, all $n\ge2$ and all $f\in\Hom_\cD(X,Y[n])$, there is an epimorphism $X'\twoheadrightarrow X$ in $\cA$ such that the composite $X'\to X\xrightarrow{f}Y[n]$ vanishes.
\item In this case $\real$ is an equivalence, since $\cD$ is generated by $\cA$ under extensions and shifts.
\end{enumerate}
\end{theorem}

\begin{remark}
In \cite[Prop.~3.1.16]{BBD82} the effaceability condition is formulated with monomorphisms $Y\hookrightarrow Y'$. The version with epimorphisms follows by applying the statement to $\cD^{\mathrm{op}}$. The existence of $\real$ for algebraic triangulated categories is established, for instance, in \cite{Let24}.
\end{remark}

\subsection{Simple-minded collections, silting objects and co-t-structures}

\begin{definition}\label{def:smc}
A finite set $\{S_1,\dots,S_r\}$ of objects of $\cD$ is a \emph{simple-minded collection} if
\begin{enumerate}
\item[(SM1)] $\Hom(S_i,S_j[l])=0$ for all $i,j$ and all $l<0$;
\item[(SM2)] $\Hom(S_i,S_j)=\delta_{ij}\kk$;
\item[(SM3)] $\thick{S_1,\dots,S_r}=\cD$.
\end{enumerate}
An object $T$ is \emph{silting} if $\Hom(T,T[l])=0$ for all $l>0$ and $\thick{T}=\cD$. It is \emph{tilting} if moreover $\Hom(T,T[l])=0$ for all $l<0$.
\end{definition}

\begin{intuition}
Simple-minded collections behave like the simple modules of a finite-dimensional algebra, and silting objects behave like its projective modules. The former determine t-structures, the latter co-t-structures (weight structures). For $\Db(\modd\Lambda)$ the two notions correspond to each other under the bijections of Koenig and Yang (\cref{thm:KY}). This ``simple versus projective'' dichotomy is the mechanism behind \cref{thm:HJY} and behind all results below.
\end{intuition}

\begin{theorem}[{\cite{AN09}; see also \cite[\S5]{KY14}}]\label{thm:smc}
Let $\{S_1,\dots,S_r\}$ be a simple-minded collection in a Hom-finite Krull--Schmidt triangulated category $\cD$. Put
\[\cD^{\le0}:=\Filt\{S_i[l]: 1\le i\le r,\ l\ge0\},\qquad \cD^{\ge0}:=\Filt\{S_i[l]:1\le i\le r,\ l\le0\}.\]
Then $(\cD^{\le0},\cD^{\ge0})$ is a bounded t-structure on $\cD$. Its heart is $\Filt\{S_1,\dots,S_r\}$, a length category whose simple objects are exactly $S_1,\dots,S_r$.
\end{theorem}

\begin{definition}[{\cite{Bon10,Pau08}}]
A \emph{co-t-structure} (or \emph{weight structure}) on $\cD$ is a pair $(\cU,\cV)$ of full subcategories with the following properties: both are closed under direct summands; $\cU[-1]\subset\cU$ and $\cV[1]\subset\cV$; $\Hom(\cU,\cV)=0$; and every object $X$ fits into a triangle $U\to X\to V\to U[1]$ with $U\in\cU$ and $V\in\cV$. Its \emph{co-heart} is $\cU\cap\cV[-1]$. It is \emph{bounded} if $\bigcup_l\cU[l]=\cD=\bigcup_l\cV[l]$.
\end{definition}

\begin{theorem}[{\cite{Bon10,AI12}}]\label{thm:silting-cot}
Let $T$ be a silting object of $\cD$. Let $\cU$ (resp.\ $\cV$) be the smallest full subcategory closed under extensions and direct summands containing all $T[l]$ with $l\le0$ (resp.\ $l\ge1$). Then $(\cU,\cV)$ is a bounded co-t-structure with co-heart $\add T$.
\end{theorem}

\begin{theorem}[{Koenig--Yang \cite[Thm.~6.1]{KY14}}]\label{thm:KY}
For a finite-dimensional algebra $\Lambda$ there are bijections between:
\begin{enumerate}
\item equivalence classes of silting objects in $\Kb(\proj\Lambda)$;
\item equivalence classes of simple-minded collections in $\Db(\modd\Lambda)$;
\item bounded t-structures on $\Db(\modd\Lambda)$ with length heart;
\item bounded co-t-structures on $\Kb(\proj\Lambda)$.
\end{enumerate}
\end{theorem}

\begin{theorem}[{Keller \cite{Kel94}}]\label{thm:keller}
Let $\cD$ be an algebraic idempotent complete triangulated category, and let $T\in\cD$ be a tilting object. Then there is a triangle equivalence $\cD\simeq\Kb(\proj\End(T))$ sending $T$ to the free module of rank one.
\end{theorem}

\subsection{Recollements and gluing}

\begin{definition}\label{def:recollement}
A full triangulated subcategory $\cN\subset\cM$ is \emph{admissible} if the inclusion $i_*\colon\cN\to\cM$ has a left adjoint $i^*$ and a right adjoint $i^!$. Put $\cQ=\cM/\cN$ (Verdier quotient), and let $j^*\colon\cM\to\cQ$ be the quotient functor. It has a left adjoint $j_!$ and a right adjoint $j_*$, both fully faithful, which identify $\cQ$ with ${}^\perp\cN$ and with $\cN^\perp$ respectively. We obtain a \emph{recollement}
\[
\begin{tikzcd}[column sep=large]
\cN \arrow[r,"i_*" description] & \cM \arrow[l,bend right=40,"i^*"'] \arrow[l,bend left=40,"i^!"] \arrow[r,"j^*" description] & \cQ \arrow[l,bend right=40,"j_!"'] \arrow[l,bend left=40,"j_*"]
\end{tikzcd}
\]
with $i^*j_!=0$, $i^!j_*=0$, $j^*i_*=0$, $j^*j_!\cong\mathrm{id}\cong j^*j_*$, and functorial triangles
\[ j_!j^*X\to X\to i_*i^*X\to,\qquad i_*i^!X\to X\to j_*j^*X\to.\]
Applying $i^!$ to the first triangle for $X=j_*Y$ gives $i^!j_!Y\cong i^*j_*Y[-1]$.
\end{definition}

\begin{theorem}[{Gluing \cite[Thm.~1.4.10 and \S1.4]{BBD82}}]\label{thm:gluing}
In a recollement as above, let $\cN$ and $\cQ$ carry t-structures. Then
\begin{align*}
\cM^{\le0}&=\{X: j^*X\in\cQ^{\le0},\ i^*X\in\cN^{\le0}\},\\
\cM^{\ge0}&=\{X: j^*X\in\cQ^{\ge0},\ i^!X\in\cN^{\ge0}\}
\end{align*}
is a t-structure on $\cM$, the \emph{glued t-structure}. With respect to it, $i_*$ and $j^*$ are t-exact, $j_!$ and $i^*$ are right t-exact, and $j_*$ and $i^!$ are left t-exact. If the given t-structures are bounded, so is the glued one.
\end{theorem}

Let $\cA_\cN$, $\cA$ and $\cA_\cQ$ be the three hearts. Write $\pH^0$ for the cohomology functors of the glued t-structure, and set ${}^p\!j_!:=\pH^0\circ j_!$ and ${}^p\!j_*:=\pH^0\circ j_*$ on $\cA_\cQ$. For $Y\in\cA_\cQ$ define the \emph{intermediate extension}
\[ j_{!*}Y:=\im\big({}^p\!j_!Y\to{}^p\!j_*Y\big). \]

\begin{proposition}[{\cite[\S1.4]{BBD82}}]\label{prop:glued-heart}
In the situation of \cref{thm:gluing}:
\begin{enumerate}
\item $i_*$ identifies $\cA_\cN$ with the Serre subcategory $\{X\in\cA:j^*X=0\}$ of $\cA$; this subcategory is closed under subobjects, quotients and extensions.
\item For $Y\in\cA_\cQ$ we have $j^*{}^p\!j_!Y\cong Y\cong j^*{}^p\!j_*Y\cong j^*j_{!*}Y$. The object ${}^p\!j_!Y$ has no non-zero quotient in $i_*\cA_\cN$, and ${}^p\!j_*Y$ has no non-zero subobject in $i_*\cA_\cN$. Hence $j_{!*}Y$ has neither. The kernel of ${}^p\!j_!Y\twoheadrightarrow j_{!*}Y$ and the cokernel of $j_{!*}Y\hookrightarrow {}^p\!j_*Y$ lie in $i_*\cA_\cN$.
\item The simple objects of $\cA$ are the $i_*S$ with $S$ simple in $\cA_\cN$, and the $j_{!*}S$ with $S$ simple in $\cA_\cQ$.
\item If $\cA_\cN$ and $\cA_\cQ$ are length categories, so is $\cA$.
\end{enumerate}
\end{proposition}

\begin{proof}[Comments on the proof]
(1) and (3) are proved in \cite[\S1.4]{BBD82} (see in particular 1.4.22--1.4.26 there, where $j_{!*}$ is introduced and the simple objects of $\cA$ are determined).

For (2): $j^*$ is t-exact, so $j^*{}^p\!j_!Y=H^0(j^*j_!Y)=Y$, and similarly for $j_*$. Moreover $\Hom_\cA({}^p\!j_!Y,i_*Z)=\Hom(j_!Y,i_*Z)=\Hom(Y,j^*i_*Z)=0$, and dually for ${}^p\!j_*Y$. The map ${}^p\!j_!Y\to{}^p\!j_*Y$ becomes an isomorphism after applying the exact functor $j^*|_\cA$, so its kernel and cokernel lie in $i_*\cA_\cN$ by~(1).

For (4): every $X\in\cA$ sits in an exact sequence
\[0\to i_*H^{-1}(i^*X)\to{}^p\!j_!j^*X\to X\to i_*H^0(i^*X)\to0,\]
obtained from the long exact cohomology sequence of $j_!j^*X\to X\to i_*i^*X$. The functor ${}^p\!j_!$ is right exact. By~(2), ${}^p\!j_!$ of a simple object has finite length, hence so does ${}^p\!j_!$ of any object of finite length. Thus $X$ has finite length.
\end{proof}

\begin{definition}\label{def:adm-filtration}
An \emph{admissible filtration} of $\cD$ is a chain $0=\cD_0\subset\cD_1\subset\dots\subset\cD_m=\cD$ of thick subcategories such that $\cD_{s-1}$ is admissible in $\cD_s$ for $1\le s\le m$.
\begin{itemize}
\item Its \emph{strata} are $\cT_s:=\cD_s/\cD_{s-1}$. We write $i_s,j_s$ for the functors of the recollement $(\cD_{s-1},\cD_s,\cT_s)$.
\item By abuse of notation we regard $j_{s!}$ and $j_{s*}$ as functors $\cT_s\to\cD$, by composing with the inclusion $\cD_s\subset\cD$.
\item Given t-structures on all $\cT_s$, the \emph{glued t-structure} on $\cD$ is defined inductively: the glued t-structure on $\cD_s$ is obtained from that on $\cD_{s-1}$ and the one on $\cT_s$ by \cref{thm:gluing}. (For $s=1$ we have $\cD_1=\cT_1$.)
\end{itemize}
By \cref{thm:gluing}, the inclusions $\cD_{s}\subset\cD$ are t-exact. Hence $\cA_s:=\cA\cap\cD_s$ is the heart of $\cD_s$, and it is a Serre subcategory of $\cA$.
\end{definition}

\subsection{Exceptional sequences}

\begin{definition}[{\cite{Bon89}; conventions of \cite{BB26}}]
An object $E$ is \emph{exceptional} if $\End^\bullet(E)=\kk$. A sequence $(E_1,\dots,E_n)$ of exceptional objects is an \emph{exceptional sequence} if $\Hom^\bullet(E_j,E_i)=0$ for $j>i$. It is \emph{full} if $\thick{E_1,\dots,E_n}=\cD$.

For an exceptional pair $(E,F)$, the left mutation $L_EF$ is defined by the triangle $\Hom^\bullet(E,F)\otimes E\to F\to L_EF$. The \emph{left dual} of $(E_1,\dots,E_n)$ is $(F_n,\dots,F_1)$, where $F_1=E_1$ and $F_i=L_{E_1}\cdots L_{E_{i-1}}E_i$.
\end{definition}

A full exceptional sequence gives the admissible filtration $\cD_s=\thick{E_1,\dots,E_s}$, with all strata equivalent to $\Db(\kk)$. In the notation of \cref{def:adm-filtration} one has $E_s=j_{s!}\kk$ and $F_s=j_{s*}\kk$ \cite[\S2.4]{BB26}. If the sequence is full, the left dual is characterised by the \emph{Hom-duality}
\begin{equation}\label{eq:hom-duality-exc}
\Hom^\bullet(E_i,F_j)\cong\delta_{ij}\kk\quad\text{(concentrated in degree $0$)}.
\end{equation}
The \emph{t-structure glued along $(E_1,\dots,E_n)$} is the one obtained by gluing the standard t-structures on $\Db(\kk)$.

\begin{theorem}[{\cite[Lemma~2.4]{BB26}}]\label{thm:BB-2.4}
Let $(E_1,\dots,E_n)$ be a full exceptional sequence with left dual $(F_n,\dots,F_1)$. The t-structure glued along it is
\[\cD^{\le0}=\{X:\Hom(X,F_s[l])=0\ \forall s,\ \forall l<0\},\qquad \cD^{\ge0}=\{X:\Hom(E_s,X[l])=0\ \forall s,\ \forall l<0\}.\]
\end{theorem}

\subsection{Highest weight categories and the Bodzenta--Bondal theorems}

\begin{definition}[{\cite{CPS88}; cf.\ \cite[\S2.1]{BB26}}]\label{def:hw}
A $\kk$-linear abelian category is \emph{Deligne finite} if it is a Hom-finite length category with finitely many simple objects and a projective generator; equivalently, it is equivalent to $\modd\Lambda$ for a finite-dimensional algebra $\Lambda$.

Let $\cA$ be Deligne finite, with simple objects $L(1),\dots,L(m)$ indexed by the totally ordered set $\{1<\dots<m\}$, projective covers $P(i)$ and injective hulls $I(i)$. Then $\cA$ is \emph{highest weight} if there are objects $\Delta(i)$ such that:
\begin{itemize}
\item[(st1)] there is an epimorphism $\Delta(i)\twoheadrightarrow L(i)$ whose kernel has composition factors $L(j)$ with $j<i$ only;
\item[(st2)] there is an epimorphism $P(i)\twoheadrightarrow\Delta(i)$ whose kernel lies in $\cF(\Delta(j):j>i)$.
\end{itemize}
The \emph{costandard objects} $\nabla(i)$ are defined dually, using injective hulls.
\end{definition}

By \cite{CPS88}, highest weight categories are exactly the module categories of quasi-hereditary algebras. \cite{BB26} also allows partial orders. Since every refinement of the partial order gives an equivalent highest weight structure \cite[Rem.~2.3]{BB26}, we restrict to total orders throughout.

\begin{definition}[{\cite{DR92}; \cite[\S2.2]{BB26}}]\label{def:univ-ext}
Let $Q,T$ be objects of an abelian category with $V:=\Ext^1(Q,T)$ finite-dimensional. The \emph{universal extension} of $Q$ by $T$ is the extension
\[0\to T\otimes V^\vee\to R\to Q\to0\]
classified by the canonical element of $\Ext^1(Q,T\otimes V^\vee)\cong V\otimes V^\vee\cong\End(V)$.

The \emph{iterated universal extension} $G=\bigoplus_{s=1}^mG_s$ of a sequence $(T_1,\dots,T_m)$ is defined recursively. For $m=1$ put $G_1=T_1$. For $m>1$, let $G'=\bigoplus_{s<m}G'_s$ be the iterated universal extension of $(T_1,\dots,T_{m-1})$. Then $G_s$ is the universal extension of $G'_s$ by $T_m$ for $s<m$, and $G_m:=T_m$.

A sequence $(E_1,\dots,E_m)$ of objects of $\cA$ is \emph{standardizable} if all $\Hom(E_i,E_j)$ and $\Ext^1(E_i,E_j)$ are finite-dimensional, and $\rad(E_i,E_j)=0=\Ext^1(E_i,E_j)$ for $i\ge j$.
\end{definition}

\begin{theorem}[{Dlab--Ringel \cite[Thm.~2]{DR92}; \cite[Thm.~2.1]{BB26}}]\label{thm:DR}
Let $(E_1,\dots,E_m)$ be a standardizable sequence in a $\kk$-linear abelian category $\cA$, with iterated universal extension $G$. Then $\widetilde\cA=\modd\End_\cA(G)$, with the order $1<\dots<m$, is a highest weight category with standard objects $\Hom_\cA(G,E_i)$. Moreover $\cF(E_1,\dots,E_m)\subset\cA$ and $\cF(\Delta(1),\dots,\Delta(m))\subset\widetilde\cA$ are equivalent.
\end{theorem}

\begin{theorem}[{\cite[Thm.~3.1]{BB26}}]\label{thm:BB-3.1}
Let $(E_1,\dots,E_n)$ be a full exceptional sequence in $\cD$, let $(F_n,\dots,F_1)$ be its left dual, and let $\cA$ be the heart of the t-structure glued along $(E_1,\dots,E_n)$.
\begin{enumerate}
\item $\cA$ is highest weight with standard objects $\{E_i\}$, costandard objects $\{F_i\}$ and the order induced by the sequence if and only if
\[\Hom(\textstyle\bigoplus_iE_i,\bigoplus_iE_i[l])=0\quad\text{and}\quad\Hom(\bigoplus_iF_i,\bigoplus_iF_i[l])=0\qquad\text{for all }l<0.\]
\item If $\cD$ is idempotent complete and dg enhanced and the conditions of \textup{(1)} hold, then $\cD\simeq\Db(\cA)$.
\end{enumerate}
\end{theorem}

\begin{theorem}[{\cite[Thm.~3.4]{BB26}}]\label{thm:BB-3.4}
Let $\cD$ carry a t-structure with heart $\cA$, and let $(E_1,\dots,E_n)$ be an exceptional sequence in $\cD$ with left dual $(F_n,\dots,F_1)$.
\begin{enumerate}
\item If $E_i\in\cD^{\le0}$ and $F_i\in\cD^{\ge0}$ for all $i$, then the t-structure of $\cD$ restricts to the t-structure on $\cC:=\thick{E_1,\dots,E_n}$ glued along $(E_1,\dots,E_n)$. Its heart is $\cA\cap\cC$.
\item If $E_i,F_i\in\cA$ for all $i$, then $\cA\cap\cC$ is highest weight with standard objects $\{E_i\}$ and costandard objects $\{F_i\}$.
\end{enumerate}
\end{theorem}

\begin{theorem}[{\cite[Thm.~3.8]{BB26}}]\label{thm:BB-3.8}
Let $\cA$ be an abelian category. Equivalence classes of highest weight structures on $\cA$ are in bijection with equivalence classes of full exceptional sequences in $\Db(\cA)$ such that both the sequence and its left dual consist of objects of $\cA$.
\end{theorem}

\begin{example}[{\cite[Ex.~5.11]{BB26}}]\label{ex:BB-5.11}
Let $f\colon X\to Y$ be the minimal resolution of an $A_1$ surface singularity, with exceptional curve $E$. Let $\cC_f=\{F\in\Db(X):Rf_*F=0\}$, and let $\mathcal{A}_f=\cC_f\cap\mathrm{Coh}(X)$.
\begin{itemize}
\item $\mathcal{A}_f$ is equivalent to $\modd\kk$, with additive generator $\cO_E(-1)$.
\item $\cO_E(-1)$ is $2$-spherical. In particular $\Ext^2(\cO_E(-1),\cO_E(-1))\cong\kk$, so $\Db(\mathcal{A}_f)\not\simeq\cC_f$.
\end{itemize}
In the terminology of \cref{sec:P-objects}, $\cO_E(-1)$ is a $\Pn12$-object generating $\cC_f$. This example already exhibits the phenomenon of \cref{prop:heart-P}(1).
\end{example}

\begin{remark}
Kalck's algebra \cite[Prop.~5.6]{Kal17} has a full exceptional sequence of modules but is not quasi-hereditary. Bodzenta and Bondal show that its standard t-structure is nevertheless glued along this sequence, and that the dual sequence does not consist of modules \cite[\S4.1]{BB26}. The condition on the dual sequence in \cref{thm:BB-3.1,thm:BB-3.8} is therefore essential.
\end{remark}

\subsection{Stratified algebras}

\begin{definition}[{\cite{CPS96,Dla96,AHLU00}}]\label{def:strat}
Let $\cA$ be Deligne finite, with simple objects $L(1),\dots,L(m)$ indexed by $\{1<\dots<m\}$, projective covers $P(i)$ and injective hulls $I(i)$. Define the following objects.
\begin{itemize}
\item The \emph{standard object} $\Delta(i)=P(i)/U_i$, where $U_i=\sum_{j>i}\tr_{P(j)}P(i)$. It is the largest quotient of $P(i)$ all of whose composition factors are $L(j)$ with $j\le i$.
\item The \emph{proper standard object} $\Deltabar(i)$: the largest quotient of $P(i)$ such that $[\Deltabar(i):L(i)]=1$ and all other composition factors are $L(j)$ with $j<i$.
\item The \emph{costandard object} $\nabla(i)$: the largest subobject of $I(i)$ all of whose composition factors are $L(j)$ with $j\le i$.
\item The \emph{proper costandard object} $\nablabar(i)$: the largest subobject of $I(i)$ such that $[\nablabar(i):L(i)]=1$ and all other composition factors are $L(j)$ with $j<i$.
\end{itemize}
We put $B_i:=\End_\cA(\Delta(i))$. The category $\cA$ (or an algebra $\Lambda$ with $\cA\simeq\modd\Lambda$) is
\begin{itemize}
\item \emph{standardly stratified} if $U_i\in\cF(\Delta(j):j>i)$ for every $i$, i.e.\ $P(i)$ has a $\Delta$-filtration with top factor $\Delta(i)$ and remaining factors $\Delta(j)$, $j>i$;
\item \emph{properly stratified} if moreover every $P(i)$ lies in $\cF(\Deltabar)$;
\item \emph{quasi-hereditary} if it is standardly stratified and $\Delta(i)=\Deltabar(i)$ for all $i$.
\end{itemize}
\end{definition}

\begin{remark}
Several inequivalent notions of ``stratified'' are in use: some allow partial preorders, and the notion of \cite{CPS96} is homological. For total orders, the definition above is equivalent to $\Lambda_\Lambda\in\cF(\Delta)$ \cite{Dla96,AHLU00}; Mazorchuk \cite{Maz04} calls this \emph{strongly standardly stratified}. Quasi-hereditary in the above sense is the same as highest weight in the sense of \cref{def:hw}.
\end{remark}

\begin{intuition}
The local algebra $B_i$ is the ``stratum algebra'' at $i$. Quasi-heredity means $B_i=\kk$ for all $i$. In the derived picture, $\Delta(i)$ plays the role of a projective generator of the $i$-th stratum, and $\nablabar(i)$ the role of its simple object (\cref{lem:strat-basic}). For $B_i\ne\kk$ these are genuinely different objects.
\end{intuition}

\begin{lemma}\label{lem:strat-basic}
Let $\cA$ be standardly stratified.
\begin{enumerate}
\item Let $X\in\cF(\Delta(j):j>i)$, and let $Y\in\cA$ have all composition factors in $\{L(j):j\le i\}$. Then $\Ext^l_\cA(X,Y)=0$ for all $l\ge0$.
\item $\Ext^l(\Delta(j),\Delta(i))=0$ for $j>i$ and all $l\ge0$, and $\Ext^l(\Delta(i),\Delta(i))=0$ for $l\ge1$. The algebra $B_i$ is local, and $\dim B_i=[\Delta(i):L(i)]$. In particular $\cA$ is quasi-hereditary if and only if $B_i\cong\kk$ for all $i$.
\item $\Ext^l(\Delta(j),\nablabar(i))=0$ for $j>i$ and all $l\ge0$. Furthermore $\Hom(\Delta(i),\nablabar(i))\cong\kk$ and $\Ext^l(\Delta(i),\nablabar(i))=0$ for $l\ge1$, while $\Hom(\Delta(j),\nablabar(i))=0$ for $j<i$.
\item Every $\Delta(i)$ has finite projective dimension, and $\Kb(\proj\cA)=\thick{\Delta(1),\dots,\Delta(m)}$.
\end{enumerate}
\end{lemma}

\begin{proof}
(1) It suffices to treat $X=\Delta(j)$ with $j>i$; we use descending induction on $j$.
\begin{itemize}
\item Since $\Delta(j)$ has top $L(j)$ and $Y$ has no composition factor $L(j)$, we get $\Hom(\Delta(j),Y)=0$.
\item The exact sequence $0\to U_j\to P(j)\to\Delta(j)\to0$, together with $\Hom(P(j),Y)=0$, gives $\Ext^1(\Delta(j),Y)\cong\Hom(U_j,Y)$ and $\Ext^{l+1}(\Delta(j),Y)\cong\Ext^l(U_j,Y)$ for $l\ge1$.
\item Since $U_j\in\cF(\Delta(t):t>j)$, the induction hypothesis gives $\Ext^\bullet(U_j,Y)=0$. For $j=m$ we have $U_m=0$.
\end{itemize}

(2) The first two vanishing statements follow from~(1) applied to $Y=\Delta(i)$, using the sequence $0\to U_i\to P(i)\to\Delta(i)\to0$ as in~(1). Since $\Delta(i)$ has simple top, it is indecomposable, so $B_i$ is local. Every map $P(i)\to\Delta(i)$ vanishes on $U_i$, because $\Hom(U_i,\Delta(i))=0$ by~(1). Hence $B_i=\Hom(P(i),\Delta(i))$, whose dimension is $[\Delta(i):L(i)]$. The last statement follows, since $\Delta(i)=\Deltabar(i)$ if and only if $[\Delta(i):L(i)]=1$.

(3) The first statement is~(1) with $Y=\nablabar(i)$.
\begin{itemize}
\item The image of a non-zero map $\Delta(j)\to\nablabar(i)$ is a quotient of $\Delta(j)$ and contains $\soc\nablabar(i)=L(i)$. This is impossible for $j<i$, and forces $\dim\Hom(\Delta(i),\nablabar(i))\le\dim\Hom(P(i),\nablabar(i))=[\nablabar(i):L(i)]=1$ for $j=i$.
\item The composite $\Delta(i)\twoheadrightarrow L(i)\hookrightarrow\nablabar(i)$ is non-zero, so $\Hom(\Delta(i),\nablabar(i))\cong\kk$.
\item The vanishing of $\Ext^{l}(\Delta(i),\nablabar(i))$ for $l\ge1$ follows from $\Ext^\bullet(U_i,\nablabar(i))=0$, which holds by~(1).
\end{itemize}

(4) We use descending induction: $\Delta(m)=P(m)$, and $0\to U_i\to P(i)\to\Delta(i)\to0$ with $U_i\in\cF(\Delta(j):j>i)$. This shows that $\pd\Delta(i)<\infty$ and that $P(i)\in\thick{\Delta(i),\dots,\Delta(m)}$.
\end{proof}

\begin{lemma}[Inductive characterisation]\label{lem:strat-induction}
Let $\cA$ be Deligne finite, with simple objects $L(1),\dots,L(m)$ and projective covers $P(1),\dots,P(m)$. Let $\cA'\subset\cA$ be the Serre subcategory of objects without composition factor $L(m)$. For $s<m$ put
\[U_s:=\tr_{P(m)}P(s)\qquad\text{and}\qquad P'(s):=P(s)/U_s,\]
so that $P'(s)$ is a projective cover of $L(s)$ in $\cA'$. Then $\cA$ is standardly stratified (for $1<\dots<m$) if and only if both of the following hold:
\begin{enumerate}
\item[(a)] for every $s<m$, the object $U_s$ is projective, i.e.\ $U_s\cong P(m)^{\oplus q_s}$ for some $q_s\ge0$;
\item[(b)] $\cA'$ is standardly stratified (for $1<\dots<m-1$).
\end{enumerate}
In this case $\Delta(m)=P(m)$. For $s<m$, the objects $\Delta(s)$, $\Deltabar(s)$ and $\nablabar(s)$ of $\cA$ coincide with those of $\cA'$.
\end{lemma}

\begin{proof}
\emph{Comparison of $\cA$ and $\cA'$.} Let $s<m$. The objects $\Delta(s)$ and $\Deltabar(s)$ are quotients of $P(s)$ with composition factors $L(t)$, $t\le s<m$. They therefore lie in $\cA'$, are quotients of $P'(s)$, and agree with the corresponding objects of $\cA'$. Similarly, $\nablabar(s)$ lies in the largest subobject of $I(s)$ belonging to $\cA'$, which is an injective hull of $L(s)$ in $\cA'$. Finally $\Delta(m)=P(m)$, since all composition factors are $\le m$.

\emph{(a) and (b) imply standard stratification.} Let $s<m$. By (b), the kernel of $P'(s)\twoheadrightarrow\Delta(s)$ lies in $\cF(\Delta(t):s<t<m)$. The kernel of the composite $P(s)\twoheadrightarrow P'(s)\twoheadrightarrow\Delta(s)$ is an extension of this kernel by $U_s\cong\Delta(m)^{\oplus q_s}$. Hence it lies in $\cF(\Delta(t):t>s)$. For $s=m$ there is nothing to show.

\emph{Standard stratification implies (a) and (b).} We first show the following claim: every $X\in\cF(\Delta)$ has a subobject $M\in\add P(m)$ with $X/M\in\cF(\Delta(t):t<m)$. We use induction on the length of a $\Delta$-filtration. Take an exact sequence $0\to X_1\to X\to\Delta(t)\to0$, and choose $M_1\subset X_1$ for $X_1$ by induction.
\begin{itemize}
\item If $t<m$, take $M=M_1$.
\item If $t=m$, the extension $0\to X_1/M_1\to X/M_1\to\Delta(m)\to0$ splits, because $\Delta(m)=P(m)$ is projective. Let $M$ be the preimage in $X$ of the summand $\Delta(m)$. Then $0\to M_1\to M\to P(m)\to0$ also splits, so $M\in\add P(m)$, and $X/M\cong X_1/M_1$.
\end{itemize}
Now let $s<m$, and let $K_s$ be the kernel of $P(s)\twoheadrightarrow\Delta(s)$, so $K_s\in\cF(\Delta(t):t>s)$. The claim gives $M\subset K_s$ with $M\in\add P(m)$ and $K_s/M\in\cF(\Delta(t):s<t<m)$.
\begin{itemize}
\item $P(s)/M$ is an extension of $\Delta(s)$ by $K_s/M$, so it lies in $\cA'$. Hence $M\supset U_s$.
\item $M$ is generated by $P(m)$, so $M\subset U_s$.
\end{itemize}
Thus $U_s=M$, which proves (a). Moreover the kernel of $P'(s)\twoheadrightarrow\Delta(s)$ is $K_s/M\in\cF(\Delta(t):s<t<m)$, which proves (b).
\end{proof}

\begin{remark}[Heredity chains]\label{rem:heredity}
Let $\cA=\modd\Lambda$, and let $e$ be a primitive idempotent with $e\Lambda=P(m)$. Then $\bigoplus_sU_s$ corresponds to the trace ideal $J=\Lambda e\Lambda$, and $\cA'\simeq\modd\Lambda/J$. Condition (a) says that $J$ is projective as a right $\Lambda$-module. If in addition $e\Lambda e=\End(P(m))=\kk$, this is the definition of a heredity ideal, and \cref{lem:strat-induction} becomes the characterisation of quasi-hereditary algebras by heredity chains \cite{CPS88,DR92}.

So standardly stratified algebras are those admitting a chain of idempotent ideals, each projective as a right module over the corresponding quotient; the local algebras $e\Lambda e$ play the role of the stratum algebras (cf.\ \cite{Dla96,CPS96}). All inductive arguments below follow this pattern: they peel off the maximal simple object, or equivalently the last stratum.
\end{remark}



% ===============================================================
\section{Local-exceptional objects and local strata}\label{sec:local}
% ===============================================================

\subsection{Local-exceptional objects}

\begin{definition}\label{def:local-exc}
An object $E$ of $\cD$ is \emph{local-exceptional} if $\Hom(E,E[l])=0$ for all $l\ne0$ and $\End(E)$ is a local algebra. A \emph{local-exceptional sequence} is a sequence $(E_1,\dots,E_m)$ of local-exceptional objects with $\Hom^\bullet(E_t,E_s)=0$ for $t>s$. It is \emph{full} if $\thick{E_1,\dots,E_m}=\cD$.
\end{definition}

\begin{example}\label{ex:local-exc}
\leavevmode
\begin{enumerate}
\item Exceptional objects are exactly the local-exceptional objects with $\End(E)=\kk$.
\item If $B$ is a finite-dimensional local algebra, then $B$ is a local-exceptional object of $\Db(\modd B)$ and of $\Kb(\proj B)$. More generally, a module $M$ over a finite-dimensional algebra $\Lambda$ is local-exceptional in $\Db(\modd\Lambda)$ if $\End(M)$ is local and $\Ext^l(M,M)=0$ for all $l\ge1$.
\item The standard modules of a standardly stratified algebra form a full local-exceptional sequence in $\Kb(\proj\Lambda)$ (\cref{thm:C}).
\item $\Pn n0$-like objects (\cref{def:Pnk}) are exactly the local-exceptional objects with endomorphism algebra $\kk[t]/(t^{n+1})$ (\cref{sec:P-zero}).
\end{enumerate}
\end{example}

\begin{intuition}
For an exceptional object $E$ we have $\thick E\simeq\Db(\kk)$, and $E$ is simultaneously the projective object and the simple object of the heart $\modd\kk$. For a local-exceptional object these two roles separate.
\begin{itemize}
\item $E$ behaves like the free module over $B=\End(E)$ (\cref{prop:local-exc}).
\item The ``simple'' object that $E$ should be paired with is not contained in $\thick E$.
\end{itemize}
Accordingly, the strata of the filtrations considered in \cref{sec:gluing,sec:criterion} are not generated by the objects $E_s$. Each of them contains both a projective-like and a simple-like object.
\end{intuition}

\begin{lemma}\label{lem:euler}
Let $B$ be a finite-dimensional local algebra, and let $Y\in\Kb(\proj B)$. Then $\chi(Y):=\sum_l(-1)^l\dim H^l(Y)$ is divisible by $\dim B$.
\end{lemma}

\begin{proof}
Finitely generated projective modules over a local algebra are free. So $Y$ is isomorphic to a bounded complex with terms $B^{r_l}$, and $\chi(Y)=\sum_l(-1)^lr_l\dim B$.
\end{proof}

\begin{proposition}\label{prop:local-exc}
Let $\cD$ be algebraic, let $E\in\cD$ be local-exceptional, and put $B=\End(E)$.
\begin{enumerate}
\item $\thick E\simeq\Kb(\proj B)$, with $E$ corresponding to $B$.
\item $\thick E$ carries the bounded co-t-structure with co-heart $\add E$.
\item If $B\ne\kk$, there is no $X\in\thick E$ with $\dim\Hom^\bullet(E,X)=1$. In particular, the standard t-structure of $\Db(\modd B)$ does not restrict to $\Kb(\proj B)$.
\end{enumerate}
\end{proposition}

\begin{proof}
(1) The object $E$ is a tilting object of $\thick E$, so this is \cref{thm:keller}. (2) follows from \cref{thm:silting-cot}.

(3) Under the equivalence of (1), $\Hom^l(E,X)\cong H^l(Y)$, where $Y$ is the perfect complex corresponding to $X$. If $\dim\Hom^\bullet(E,X)=1$, then $\chi(Y)=\pm1$. This is impossible by \cref{lem:euler}, since $\dim B\ge2$.

For the last assertion, let $P_1\to P_0\to\kk\to0$ be a minimal projective presentation of the simple $B$-module, and let $Y=(P_1\to P_0)$ in degrees $-1,0$. If the standard t-structure restricted to $\Kb(\proj B)$, then $H^0(Y)\cong\kk$ would be perfect. But $\Hom^\bullet(B,\kk)\cong\kk$, contradicting what we just proved.
\end{proof}

\begin{remark}\label{rem:local-exc-no-t}
For $B=\kk[t]/(t^{n+1})$ with $n\ge1$, the category $\thick E$ has no non-trivial t-structure at all. This is due to Bonfert \cite[Cor.~3.10]{Bonf26} (\cref{thm:Bonfert}(3)), and for $n=1$ to Holm, J\o rgensen and Yang (\cref{thm:HJY}). Part~(3) above is a weak form of this, valid for every local algebra $B$. It is the form needed in \cref{thm:C}.
\end{remark}

\subsection{Local strata}

\begin{definition}\label{def:local-stratum}
A \emph{local stratum} is a triangulated category $\cT$ (Hom-finite and idempotent complete) together with a bounded t-structure whose heart $\cH$ is a length category with a unique simple object $S$ up to isomorphism.

The local stratum is \emph{of projective type} if $S$ has a projective cover $\Pi$ in $\cH$. In that case $B:=\End(\Pi)$ is a finite-dimensional local algebra, and $\Hom(\Pi,-)$ induces an equivalence $\cH\simeq\modd B$.
\end{definition}

\begin{example}\label{ex:local-strata}
\leavevmode
\begin{enumerate}
\item $\Db(\modd B)$ with its standard t-structure, for a finite-dimensional local algebra $B$, is of projective type, with $\Pi=B$ and $S=\kk$. The case $B=\kk$ is the stratum of an exceptional object.
\item For a $\Pn nk$-like object $P$ with $k\ge1$, the category $\thick P$ with the t-structure of \cref{prop:tP} is a local stratum with $S=P$. It is of projective type if and only if $n=0$, or $k\ge2$, or ($k=1$ and $n\ge2$) (\cref{prop:heart-P}).
\end{enumerate}
\end{example}

\begin{intuition}
In a local stratum of projective type, the two roles played by $\kk$ in $\Db(\kk)$ are taken by $\Pi$ (projective-like) and $S$ (simple-like). They are related by $\Hom(\Pi,S)\cong\kk$. The remaining groups $\Hom(\Pi,S[l])$ with $l\ge1$ measure how far $\cT$ is from $\Db(\modd B)$ (\cref{lem:local-stratum-derived}).
\end{intuition}

The following lemma describes the aisles of a local stratum in terms of two ``test objects''. It is the local ingredient in the description of glued aisles (\cref{prop:aisles}).


\begin{lemma}\label{lem:aisles}
Let $\cT$ carry a bounded t-structure whose heart $\cH$ is a length category with a unique simple object $S$ (up to isomorphism). Let $Q\in\cH$ satisfy $\Hom(Q,M)\ne0$ for every non-zero $M\in\cH$. Then
\[\cT^{\le0}=\{X:\Hom(X,S[l])=0\ \forall l<0\},\qquad\cT^{\ge0}=\{X:\Hom(Q,X[l])=0\ \forall l<0\}.\]
Moreover $\cT^{\le0}=\Filt\{S[l]:l\ge0\}$ and $\cT^{\ge0}=\Filt\{S[l]:l\le0\}$. The hypothesis on $Q$ holds for $Q=S$, and for a projective cover of $S$ if one exists.
\end{lemma}

\begin{proof}
\textbf{Description of $\cT^{\ge0}$.} If $X\in\cT^{\ge0}$ and $l<0$, then $X[l]\in\cT^{\ge1}$, so $\Hom(Q,X[l])=0$. Conversely, suppose $X\notin\cT^{\ge0}$. By boundedness there is a minimal $d<0$ with $H^d(X)\ne0$, and we have a triangle $H^d(X)[-d]\to X\to\tau^{\ge d+1}X\to$. Shifting by $d$ and applying $\Hom(Q,-)$ gives an exact sequence
\[\Hom(Q,(\tau^{\ge d+1}X)[d-1])\to\Hom(Q,H^d(X))\to\Hom(Q,X[d]).\]
The first term vanishes because $(\tau^{\ge d+1}X)[d-1]\in\cT^{\ge2}$. The middle term is non-zero by hypothesis. Hence $\Hom(Q,X[d])\ne0$ with $d<0$.

\textbf{Description of $\cT^{\le0}$.} If $X\in\cT^{\le0}$ and $l<0$, then $S[l]\in\cT^{\ge1}$, so $\Hom(X,S[l])=0$. Conversely, suppose $X\notin\cT^{\le0}$, and let $d>0$ be maximal with $H^d(X)\ne0$. Since $X\in\cT^{\le d}$, for every $M\in\cH$ we have
\[\Hom(X,M[-d])\cong\Hom(\tau^{\ge d}X,M[-d])\cong\Hom(H^d(X),M).\]
Taking $M=S$ and using an epimorphism $H^d(X)\twoheadrightarrow S$, we obtain $\Hom(X,S[-d])\ne0$ with $-d<0$.

\textbf{The remaining statements.} The descriptions via $\Filt$ hold because $\cH=\Filt(S)$ and the t-structure is bounded. Every non-zero $M\in\cH$ has a simple subobject (isomorphic to $S$) and a simple quotient; the projection from a projective cover of $S$ onto the simple quotient lifts to $M$. This gives the last assertion.
\end{proof}

\begin{lemma}\label{lem:recognise-proj}
Let $\cT$ be a local stratum with simple object $S$, and let $E\in\cT$ satisfy:
\begin{itemize}
\item $\Hom(E,E[l])=0$ for all $l<0$;
\item $\Hom^\bullet(E,S)\cong\kk$, concentrated in degree $0$.
\end{itemize}
Then $E$ lies in the heart and is a projective cover of $S$. In particular $\cT$ is of projective type, with $\Pi=E$ and $\Hom(\Pi,S[l])=0$ for all $l\ne0$.
\end{lemma}

\begin{proof}
By \cref{lem:aisles}, $E\in\cT^{\le0}$. Since $\Hom(E,S)\cong\Hom(H^0(E),S)\neq0$, we have $H^0(E)\ne0$.

Suppose $E\notin\cT^{\ge0}$, and let $a<0$ be minimal with $H^a(E)\ne0$. We use the argument of \cite[Lemma~3.12]{Bonf26}. The triangle $(\tau^{\ge a+1}E)[a-1]\to H^a(E)\to E[a]\to$ (obtained by shifting the truncation triangle of $E$) gives an exact sequence
\[\Hom(E,(\tau^{\ge a+1}E)[a-1])\to\Hom(E,H^a(E))\to\Hom(E,E[a]).\]
The first term vanishes because $(\tau^{\ge a+1}E)[a-1]\in\cT^{\ge2}$ and $E\in\cT^{\le0}$. Since $E\in\cT^{\le0}$ and $H^a(E)$ lies in the heart, we have
\[\Hom(E,H^a(E))\cong\Hom(H^0(E),H^a(E)).\]
This is non-zero: $H^0(E)$ and $H^a(E)$ are non-zero objects of a length category with a unique simple object $S$, so there is a non-zero map $H^0(E)\twoheadrightarrow S\hookrightarrow H^a(E)$. Hence $\Hom(E,E[a])\ne0$ with $a<0$, a contradiction.

So $E$ lies in the heart. Since $\Ext^1(E,S)=\Hom(E,S[1])=0$ and $S$ is the only simple object, $E$ is projective. Since $\Hom(E,S)\cong\kk$, its top is $S$.
\end{proof}

\begin{lemma}\label{lem:local-stratum-derived}
Let $\cT$ be a local stratum of projective type, with $\Hom(\Pi,S[l])=0$ for all $l\ge1$. Then:
\begin{enumerate}
\item $\Hom(\Pi,Y[l])=0$ for all $Y\in\cH$ and $l\ge1$; in particular $\Pi$ is local-exceptional;
\item if $\cT$ is algebraic, then $\real\colon\Db(\modd B)\simeq\Db(\cH)\to\cT$ is an equivalence.
\end{enumerate}
Conversely, if $\real$ is an equivalence, then $\Hom(\Pi,S[l])=0$ for all $l\ge1$.
\end{lemma}

\begin{proof}
(1) follows by d\'evissage, since all composition factors of $Y$ are isomorphic to $S$. 

(2) Every $X\in\cH$ is a quotient of some $\Pi^{\oplus r}$. By (1), every morphism $X\to Y[n]$ with $n\ge1$ is effaced by such an epimorphism. So \cref{thm:realisation} applies.

For the converse, use $\Ext^l_B(B,\kk)=0$ for $l\geq1$. This is the case $m=1$ of \cref{thm:A}(3).
\end{proof}

% ===============================================================
\section{Gluing strata with local hearts}\label{sec:gluing}
% ===============================================================

\subsection{Setting}

\begin{setting}\label{set:local}
Let $\cD$ be a triangulated category (Hom-finite and idempotent complete, as always) with an admissible filtration $0=\cD_0\subset\cD_1\subset\dots\subset\cD_m=\cD$ (\cref{def:adm-filtration}) with strata $\cT_s=\cD_s/\cD_{s-1}$. Assume:
\begin{itemize}
\item each $\cT_s$ is a local stratum (\cref{def:local-stratum}): it carries a bounded t-structure whose heart $\cH_s$ is a length category with a unique simple object $S_s$ up to isomorphism;
\item $\cD$ carries the glued t-structure, with heart $\cA$; we write $\cA_s=\cA\cap\cD_s$ for the heart of $\cD_s$.
\end{itemize}
We say that we are in the \emph{projective case} if all strata are of projective type, i.e.\ each $S_s$ has a projective cover $\Pi_s$ in $\cH_s$. In this case we put $B_s:=\End_{\cH_s}(\Pi_s)$; this is a finite-dimensional local algebra, and $\cH_s\simeq\modd B_s$.
\end{setting}

\begin{notation}\label{not:objects}
In \cref{set:local} we define:
\begin{itemize}
\item $L_s:=j_{s!*}S_s\in\cA_s\subset\cA$; by \cref{prop:glued-heart}(3) and induction, $L_1,\dots,L_m$ are the simple objects of $\cA$;
\item $\nablabar_s:=j_{s*}S_s\in\cD_s\subset\cD$.
\end{itemize}
A family $Q=(Q_s)_s$ of objects $Q_s\in\cH_s$ with $\Hom(Q_s,M)\ne0$ for all non-zero $M\in\cH_s$ is called a family of \emph{test objects}; for instance $Q_s=S_s$, or $Q_s=\Pi_s$ in the projective case (\cref{lem:aisles}). For such a family we put $\Delta^Q_s:=j_{s!}Q_s$. In the projective case we write
\[\Delta_s:=\Delta^\Pi_s=j_{s!}\Pi_s .\]
\end{notation}

\begin{example}\label{ex:local-settings}
\leavevmode
\begin{enumerate}
\item \emph{Exceptional sequences.} If $\cT_s\simeq\Db(\kk)$ with the standard t-structure, then $\Pi_s=S_s=\kk$, and $\Delta_s=E_s$, $\nablabar_s=F_s$ are a dual pair of exceptional sequences.
\item \emph{Derived categories of local algebras.} If $\cT_s\simeq\Db(\modd B_s)$ with the standard t-structure, then $\Pi_s=B_s$ and $S_s=\kk$.
\item \emph{$\Pn nk$-like objects.} If $\cT_s=\thick{P_s}$ with $P_s$ $\Pn{n_s}{k_s}$-like, $k_s\ge1$, and $t_{P_s}$ as in \cref{prop:tP}, then $S_s=P_s$. By \cref{prop:heart-P}, we are in the projective case if and only if $n_s=0$, or $k_s\ge2$, or ($k_s=1$ and $n_s\ge2$).
\end{enumerate}
\end{example}

\begin{intuition}
$\Delta_s$ is the projective object of the $s$-th stratum extended by zero, and $\nablabar_s$ is the simple object of the $s$-th stratum extended by $j_{s*}$. In the exceptional case the projective and the simple object of the stratum coincide, so $\Delta_s$ and $\nablabar_s$ belong to a dual pair of exceptional sequences. In general they are of a different nature: one is ``projective-like'', the other ``simple-like''. This asymmetric pairing is what makes the Hom-duality of \cref{lem:hom-duality} hold. It replaces the Hom-duality \eqref{eq:hom-duality-exc} of exceptional sequences, which is not available, since mutation does not exchange projectives and simples.
\end{intuition}

\subsection{Hom-duality and the glued aisles}

\begin{lemma}[Hom-duality]\label{lem:hom-duality}
In \cref{set:local}, with a family of test objects $Q$:
\begin{enumerate}
\item $\Hom^\bullet(\Delta^Q_s,\nablabar_t)=0$ for $s\ne t$;
\item $\Hom^\bullet(\Delta^Q_s,\nablabar_s)\cong\Hom^\bullet_{\cT_s}(Q_s,S_s)$. This is concentrated in degrees $\ge0$, and its degree $0$ part is non-zero.
\end{enumerate}
In the projective case, $\Hom^0(\Delta_s,\nablabar_s)\cong\kk$.
\end{lemma}

\begin{proof}
Recall that $j_{s!}$ identifies $\cT_s$ with ${}^\perp\cD_{s-1}\cap\cD_s$, and $j_{s*}$ identifies it with $\cD_{s-1}^\perp\cap\cD_s$.
\begin{itemize}
\item If $s<t$, then $\Delta^Q_s\in\cD_s\subset\cD_{t-1}$ while $\nablabar_t\in\cD_{t-1}^\perp$.
\item If $s>t$, then $\nablabar_t\in\cD_t\subset\cD_{s-1}$ while $\Delta^Q_s\in{}^\perp\cD_{s-1}$.
\item If $s=t$, adjunction and $j_s^*j_{s!}\cong\mathrm{id}$ give $\Hom(j_{s!}Q_s,j_{s*}S_s[l])\cong\Hom(Q_s,S_s[l])$. This vanishes for $l<0$ since $Q_s,S_s\in\cH_s$, and it is non-zero for $l=0$ by the definition of test objects.
\end{itemize}
For a projective cover $\Pi_s$ we have $\dim\Hom(\Pi_s,S_s)=[S_s:S_s]=1$.
\end{proof}

\begin{proposition}[Glued aisles; cf.\ {\cref{thm:BB-2.4}}]\label{prop:aisles}
In \cref{set:local}, with a family of test objects $Q$,
\begin{align*}
\cD^{\le0}&=\{X\in\cD:\Hom(X,\nablabar_s[l])=0\ \forall s,\ \forall l<0\},\\
\cD^{\ge0}&=\{X\in\cD:\Hom(\Delta^Q_s,X[l])=0\ \forall s,\ \forall l<0\}.
\end{align*}
\end{proposition}

\begin{proof}
We use induction on $m$; the case $m=1$ is \cref{lem:aisles}. For $m\ge2$ write $i=i_m$ and $j=j_m$. By \cref{thm:gluing}, $X\in\cD^{\le0}$ if and only if $i^*X\in\cD_{m-1}^{\le0}$ and $j^*X\in\cT_m^{\le0}$.
\begin{itemize}
\item By induction, $i^*X\in\cD_{m-1}^{\le0}$ if and only if $\Hom(i^*X,\nablabar_s[l])=0$ for $s<m$ and $l<0$. By adjunction this is $\Hom(X,i_*\nablabar_s[l])$, and the objects $\nablabar_s$ ($s<m$) of $\cD_{m-1}$ and of $\cD$ agree.
\item By \cref{lem:aisles}, $j^*X\in\cT_m^{\le0}$ if and only if $\Hom(j^*X,S_m[l])=\Hom(X,\nablabar_m[l])$ vanishes for $l<0$.
\end{itemize}
The description of $\cD^{\ge0}$ is proved in the same way, using $i^!$, the adjunction $\Hom(i_*\Delta^Q_s,X)=\Hom(\Delta^Q_s,i^!X)$, and $\Hom(Q_m,j^*X[l])=\Hom(\Delta^Q_m,X[l])$. This is the argument of \cite[Lemma~2.4]{BB26}, with $\Hom^\bullet(-,F)^\vee$ and $\Hom^\bullet(E,-)$ replaced by the two different test functors of \cref{lem:aisles}.
\end{proof}

\begin{proposition}[Criterion; cf.\ {\cref{thm:BB-3.1}(1)}]\label{prop:criterion}
In \cref{set:local}, with a family of test objects $Q$, the following are equivalent:
\begin{enumerate}
\item[(a)] $\Delta^Q_s\in\cA$ and $\nablabar_s\in\cA$ for all $s$;
\item[(b)] $\Hom(\Delta^Q_s,\Delta^Q_t[l])=0$ and $\Hom(\nablabar_s,\nablabar_t[l])=0$ for all $s,t$ and all $l<0$.
\end{enumerate}
If these conditions hold, then $j_{s*}M\in\cA$ and $j_{s!}N\in\cA$ for all $M\in\cH_s$ and all $N\in\Filt(Q_s)\subset\cH_s$. In the projective case, if $j_{s!}S_s\in\cA$, then also $\Delta_s\in\cA$.
\end{proposition}

\begin{proof}
(a)$\Rightarrow$(b) holds because $\cA$ is a heart.

(b)$\Rightarrow$(a): By \cref{prop:aisles}, $\Delta^Q_s\in\cD^{\ge0}$ is the first condition in (b), and $\Delta^Q_s\in\cD^{\le0}$ follows from \cref{lem:hom-duality}. The argument for $\nablabar_s$ is dual.

The remaining assertions hold because $\cH_s=\Filt(S_s)$, $\Pi_s\in\Filt(S_s)$, and $j_{s*}$ and $j_{s!}$ are triangulated functors, while $\cA$ is closed under extensions.
\end{proof}

\begin{proposition}[Recognising a glued t-structure; cf.\ {\cref{thm:BB-3.4}(1)}]\label{prop:recognise}
Let $\cD$ be as in \cref{set:local}, and suppose $\cD$ carries a second t-structure $(\cD_\natural^{\le0},\cD_\natural^{\ge0})$. If, for some family of test objects, $\Delta^Q_s\in\cD_\natural^{\le0}$ and $\nablabar_s\in\cD_\natural^{\ge0}$ for all $s$, then $(\cD_\natural^{\le0},\cD_\natural^{\ge0})$ coincides with the glued t-structure.
\end{proposition}

\begin{proof}
Let $X\in\cD_\natural^{\le0}$ and $l<0$. Then $\nablabar_s[l]\in\cD_\natural^{\ge1}$, so $\Hom(X,\nablabar_s[l])=0$. By \cref{prop:aisles} this gives $\cD_\natural^{\le0}\subset\cD^{\le0}$. Dually, $\cD_\natural^{\ge0}\subset\cD^{\ge0}$. Two t-structures with $\cD_\natural^{\le0}\subset\cD^{\le0}$ and $\cD_\natural^{\ge0}\subset\cD^{\ge0}$ coincide. Indeed, the second inclusion gives $\cD^{\le0}={}^\perp(\cD^{\ge1})\subset{}^\perp(\cD_\natural^{\ge1})=\cD_\natural^{\le0}$, so the aisles agree, and hence so do the co-aisles.
\end{proof}

\subsection{The induction step}
The theory of quasi-hereditary algebras proceeds by induction along a heredity chain: one splits off the heredity ideal of a maximal weight and passes to the quotient. We follow the same pattern.
\begin{itemize}
\item On the abelian side, \cref{lem:strat-induction} reduces the standard stratification of $\cA$ to that of the Serre subcategory $\cA'$ of objects without the maximal simple object, together with a projectivity condition on traces.
\item On the triangulated side, the following proposition peels off the last stratum of the filtration and verifies this trace condition.
\end{itemize}
\Cref{thm:A,thm:B} are then proved by induction on the number of strata.

For a stratum of projective type we consider the conditions
\begin{align}
&\Hom_{\cT}(\Pi,S[2])=0;\label{eq:star}\tag{$\star$}\\
&\Hom_{\cT}(\Pi,S[l])=0\qquad\text{for all }l\ge1.\label{eq:starstar}\tag{$\star\star$}
\end{align}
For a filtration we require them for every stratum $\cT_s$, with $\Pi=\Pi_s$ and $S=S_s$. Since $\Hom(\Pi,S[1])=\Ext^1_\cH(\Pi,S)=0$, condition \eqref{eq:star} is the first non-trivial instance of \eqref{eq:starstar}. For uniform notation we let $r=2$ in the situation of \eqref{eq:star} and $r=\infty$ in that of \eqref{eq:starstar}, so that the condition reads $\Hom_\cT(\Pi,S[l])=0$ for $1\le l\le r$.

\begin{setting}[One step]\label{set:step}
Let $\cD'\subset\cD$ be an admissible subcategory, with recollement $(\cD',\cD,\cT)$ and functors $i,j$. Assume that:
\begin{itemize}
\item $\cT$ is a local stratum of projective type, with simple object $S$, projective cover $\Pi$ and $B=\End(\Pi)$;
\item $\cD'$ carries a bounded t-structure whose heart $\cA'$ is a length category.
\end{itemize}
Let $\cA$ be the heart of the glued t-structure on $\cD$; thus $i_*$ identifies $\cA'$ with the Serre subcategory $\{X\in\cA:j^*X=0\}$. Put
\[\Delta:=j_!\Pi,\qquad\nablabar:=j_*S,\qquad L:=j_{!*}S.\]
\end{setting}

\begin{proposition}[Induction step]\label{prop:step}
In \cref{set:step}, fix $r\in\{2,\infty\}$ and assume:
\begin{enumerate}[label=\textup{(\roman*)}]
\item $\Delta\in\cA$ and $\nablabar\in\cA$;
\item $\Hom_\cT(\Pi,S[l])=0$ for $1\le l\le r$;
\item $\cA'$ has a projective generator $G'$ with $\Hom_\cD(G',Y[l])=0$ for all $Y\in\cA'$ and $1\le l\le r$.
\end{enumerate}
Then:
\begin{enumerate}
\item $\Delta$ is a projective cover of the simple object $L$ in $\cA$, $\End_\cA(\Delta)\cong B$, and $\Hom_\cD(\Delta,Y[l])=0$ for all $Y\in\cA$ and $1\le l\le r$.
\item Write $G'=\bigoplus_sG'_s$, and let $G_s$ be the universal extension of $G'_s$ by $\Delta$. Then $G:=\bigoplus_sG_s\oplus\Delta$ is a projective generator of $\cA$ with $\Hom_\cD(G,Y[l])=0$ for all $Y\in\cA$ and $1\le l\le r$.
\item For every projective object $P$ of $\cA$, the trace $\tr_\Delta(P)$ is isomorphic to $\Delta^{\oplus q}$ for some $q$, and $P/\tr_\Delta(P)$ is a projective object of $\cA'$.
\item $\nablabar$ has simple socle $L$ and $[\nablabar:L]=1$; all other composition factors of $\nablabar$ lie in $\cA'$; and $\Ext^1_\cA(Y,\nablabar)=0$ for all $Y\in\cA'$.
\end{enumerate}
\end{proposition}

\begin{proof}
\textbf{Preliminaries.} Recall that $\Ext^1_\cA=\Hom_\cD(-,-[1])$ on $\cA$, and that $\Ext^2_\cA$ embeds into $\Hom_\cD(-,-[2])$ (\cref{prop:ext-comparison}). By \cref{prop:glued-heart}, the simple objects of $\cA$ are $L$ and the $i_*L'$ with $L'$ simple in $\cA'$, and $\cA$ is a length category. We use two exact sequences in $\cA$:
\begin{itemize}
\item Since $\nablabar\in\cA$, we have ${}^p\!j_*S=\nablabar$, and \cref{prop:glued-heart}(2) gives
\begin{equation}\label{eq:nabla-seq}
0\to L\to\nablabar\to C\to0\qquad\text{with }C\in i_*\cA'.
\end{equation}
\item Since $\Hom_\cA(\Delta,L)\cong\Hom(\Pi,j^*L)=\Hom(\Pi,S)\cong\kk$ and $L$ is simple, there is
\begin{equation}\label{eq:delta-seq}
0\to K\to\Delta\to L\to0\qquad\text{with }K\in\cA.
\end{equation}
\end{itemize}
We also use the following orthogonality relations, which come from $i^*j_!=0$, $i^!j_*=0$ and adjunction:
\begin{equation}\label{eq:step-orth}
\Hom^\bullet(\Delta,i_*Y)=0,\qquad\Hom^\bullet(i_*Y,\nablabar)=0\quad(Y\in\cA'),\qquad\Hom^\bullet(\Delta,\nablabar)\cong\Hom^\bullet_\cT(\Pi,S).
\end{equation}

\textbf{(1)} Let $1\le l\le r$. By \eqref{eq:nabla-seq} there is an exact sequence $\Hom(\Delta,C[l-1])\to\Hom(\Delta,L[l])\to\Hom(\Delta,\nablabar[l])$. Its outer terms vanish by \eqref{eq:step-orth} and (ii). Together with $\Hom(\Delta,i_*Y[l])=0$, d\'evissage gives $\Hom(\Delta,Y[l])=0$ for all $Y\in\cA$. In particular $\Delta$ is projective. By \eqref{eq:step-orth} and $\Hom(\Delta,L)\cong\kk$, its top is $L$. Finally $\End(\Delta)\cong\End_\cT(\Pi)=B$, since $j_!$ is fully faithful.

\textbf{(2)} For each $s$ put $V_s:=\Hom(G'_s,\Delta[1])=\Ext^1_\cA(G'_s,\Delta)$; this is finite-dimensional. The universal extension is a triangle
\begin{equation}\label{eq:univ-tri}
\Delta\otimes V_s^\vee\to G_s\to G'_s\xrightarrow{\ \varepsilon\ }\Delta\otimes V_s^\vee[1],
\end{equation}
where $\varepsilon$ is the canonical element. We prove the vanishing for $G_s$ in five steps; the case of $\Delta$ is (1).
\begin{enumerate}[label=(\alph*)]
\item \emph{Objects of $\cA'$.} Since $i^*\Delta=0$, we have $i^*G_s\cong G'_s$. Hence $\Hom(G_s,i_*Y[l])\cong\Hom(G'_s,Y[l])=0$ for $Y\in\cA'$ and $1\le l\le r$, by (iii).
\item \emph{$\Hom(G_s,\Delta[1])=0$.} Apply $\Hom(-,\Delta[1])$ to \eqref{eq:univ-tri}. The map $\varepsilon^*\colon\Hom(\Delta\otimes V_s^\vee,\Delta)\to\Hom(G'_s,\Delta[1])=V_s$ is surjective: on $\kk\cdot\mathrm{id}\otimes V_s$ it is the identity. Moreover $\Hom(\Delta\otimes V_s^\vee,\Delta[1])=0$ by~(1).
\item \emph{Morphisms to $\nablabar$.} We have $\Hom^\bullet(G'_s,\nablabar)=0$ by \eqref{eq:step-orth}. So \eqref{eq:univ-tri} gives $\Hom^\bullet(G_s,\nablabar)\cong V_s\otimes\Hom^\bullet_\cT(\Pi,S)$, which vanishes in degrees $1\le l\le r$ by (ii).
\item \emph{Degrees $2\le l\le r$.} By \eqref{eq:nabla-seq} there is an exact sequence $\Hom(G_s,C[l-1])\to\Hom(G_s,L[l])\to\Hom(G_s,\nablabar[l])$. Its outer terms vanish by (a) and (c). With (a) and d\'evissage, $\Hom(G_s,Y[l])=0$ for all $Y\in\cA$ and $2\le l\le r$.
\item \emph{Degree $1$.} By \eqref{eq:delta-seq} there is an exact sequence $\Hom(G_s,\Delta[1])\to\Hom(G_s,L[1])\to\Hom(G_s,K[2])$. Its outer terms vanish by (b) and (d). With (a), $\Hom(G_s,Y[1])=0$ for all $Y\in\cA$.
\end{enumerate}
So $G$ is projective. It is a generator: $\Hom(G_s,i_*L')\cong\Hom(G'_s,L')$ is non-zero for some $s$ because $G'$ is a generator of $\cA'$, and $\Hom(\Delta,L)\ne0$.

\textbf{(3)} Let $U:=\tr_\Delta(P)$ and $P':=P/U$.
\begin{itemize}
\item \emph{$P'$ is a projective object of $\cA'$.} Every morphism from $\Delta$ to $P'$ lifts to $P$, so it has image in $U$ and vanishes. Since $\Delta$ is a projective cover of $L$, this gives $[P':L]=0$, i.e.\ $P'\in\cA'$. Every morphism from $P$ to an object of $\cA'$ vanishes on $U$, by \eqref{eq:step-orth}. Hence $\Hom(P',-)=\Hom(P,-)$ on $\cA'$, and $P'$ is projective in $\cA'$.
\item \emph{$j^*U\cong\Pi^{\oplus q}$.} Since $j^*P'=0$, we have $j^*U=j^*P$. For $M\in\cH$ we have $j_*M\in\Filt(\nablabar)\subset\cA$, so
\[\Ext^1_\cH(j^*P,M)=\Hom(P,j_*M[1])=\Ext^1_\cA(P,j_*M)=0.\]
Hence $j^*U$ is projective in $\cH$, i.e.\ $j^*U\cong\Pi^{\oplus q}$.
\item \emph{The counit is an epimorphism.} The counit $c\colon j_!j^*U\cong\Delta^{\oplus q}\to U$ is a morphism in $\cA$ with $j^*c$ an isomorphism. By adjunction every morphism $\Delta\to U$ factors through $c$, and $U$ is the sum of the images of such morphisms; so $c$ is an epimorphism.
\item \emph{The counit is a monomorphism.} Let $N:=\ker c$. Then $j^*N=0$, so $N\in i_*\cA'$. For $Y\in\cA'$, the sequences $0\to N\to\Delta^{\oplus q}\to U\to0$ and $0\to U\to P\to P'\to0$, together with \eqref{eq:step-orth} and the projectivity of $P$, give
\[\Hom(N,i_*Y)\cong\Ext^1_\cA(U,i_*Y)\cong\Ext^2_\cA(P',i_*Y)\hookrightarrow\Hom_\cD(P',Y[2]).\]
The last group vanishes by (iii), since $P'\in\add G'$ and $r\ge2$. Taking $i_*Y=N$ gives $N=0$.
\end{itemize}

\textbf{(4)} The functor $j^*$ is exact on $\cA$, kills $\cA'$ and sends $L$ to $S$. So $[\nablabar:L]$ is the length of $j^*\nablabar=S$, which is $1$. By \eqref{eq:step-orth}, $\Hom(i_*Y,\nablabar)=0$ and $\Ext^1(i_*Y,\nablabar)=\Hom(i_*Y,\nablabar[1])=0$ for $Y\in\cA'$. Hence the socle of $\nablabar$ contains no composition factor from $\cA'$, so it is $L$.
\end{proof}

\subsection{The main theorem}
\begin{theorem}\label{thm:A}
Assume the projective case of \cref{set:local}, that $\Delta_s,\nablabar_s\in\cA$ for all $s$ (equivalently, condition \textup{(b)} of \cref{prop:criterion}), and that \eqref{eq:star} holds. Let $G=\bigoplus_{s=1}^mG_s$ be the iterated universal extension (\cref{def:univ-ext}) of $(\Delta_1,\dots,\Delta_m)$ in $\cA$.
\begin{enumerate}
\item $G$ is a projective generator of $\cA$, and $\Hom_\cD(G,Y[2])=0$ for all $Y\in\cA$. Hence $\cA\simeq\modd\Lambda$ with $\Lambda=\End_\cA(G)$.
\item With respect to the order $1<\dots<m$ of its simple objects $L_1,\dots,L_m$, the category $\cA$ is standardly stratified, with:
\begin{itemize}
\item standard objects $\Delta(s)\cong\Delta_s$, and $\End_\cA(\Delta_s)\cong B_s$;
\item proper costandard objects $\nablabar(s)\cong\nablabar_s$.
\end{itemize}
It is quasi-hereditary if and only if $B_s\cong\kk$ for all $s$.
\item Suppose $\cD$ is algebraic. Then $\real\colon\Db(\cA)\to\cD$ is an equivalence if and only if \eqref{eq:starstar} holds. In this case:
\begin{itemize}
\item $G$ is a tilting object of $\thick{\Delta_1,\dots,\Delta_m}$, and $\real$ identifies $\Kb(\proj\cA)$ with $\thick{\Delta_1,\dots,\Delta_m}$;
\item $\thick{\Delta_1,\dots,\Delta_m}=\cD$ if and only if $\cA$ is quasi-hereditary;
\item each stratum satisfies $\cT_s\simeq\Db(\modd B_s)$.
\end{itemize}
\end{enumerate}
\end{theorem}

\begin{proof}
Put $r=2$, and $r=\infty$ when proving sufficiency in (3). We prove by induction on $m$ the following statement.
\begin{quote}
$(\mathrm I_m)$: $\cA$ is standardly stratified with respect to $1<\dots<m$, with $\Delta(s)\cong\Delta_s$, $\nablabar(s)\cong\nablabar_s$ and $\End(\Delta(s))\cong B_s$. Moreover, the iterated universal extension $G^{(m)}$ of $(\Delta_1,\dots,\Delta_m)$ is a projective generator of $\cA$ with $\Hom(G^{(m)},Y[l])=0$ for all $Y\in\cA$ and $1\le l\le r$.
\end{quote}

\emph{Base case $m=1$.} Here $\cA=\cH_1\simeq\modd B_1$, which has a single simple object. Hence it is standardly stratified, with $\Delta(1)=\Pi_1$ and $\nablabar(1)=S_1$; note that $j_1=\mathrm{id}$. We have $G^{(1)}=\Pi_1$, and $\Hom(\Pi_1,Y[l])=0$ for $1\le l\le r$ by the assumption on $\cT_1$ and d\'evissage.

\emph{Inductive step $m\ge2$.} The filtration $\cD_1\subset\dots\subset\cD_{m-1}$ of $\cD_{m-1}$ satisfies all the hypotheses:
\begin{itemize}
\item it has the same strata;
\item its glued heart is $\cA_{m-1}=\cA\cap\cD_{m-1}$;
\item its objects $\Delta_s,\nablabar_s$ (for $s<m$) are the same, and lie in $\cA_{m-1}$.
\end{itemize}
So $(\mathrm I_{m-1})$ holds for $\cA_{m-1}$. Apply \cref{prop:step} to the recollement $(\cD_{m-1},\cD,\cT_m)$, with $G'=G^{(m-1)}$. Hypothesis (i) is assumed, (ii) is the assumption on $\cT_m$, and (iii) is $(\mathrm I_{m-1})$.
\begin{itemize}
\item By \cref{prop:step}(2), $G^{(m)}=G$ is a projective generator with the required vanishing. In particular $\cA$ is Deligne finite.
\item By \cref{prop:step}(1), $P(m)=\Delta_m$ with $\End(\Delta_m)\cong B_m$. By \cref{prop:step}(3), condition (a) of \cref{lem:strat-induction} holds. Condition (b) is $(\mathrm I_{m-1})$. Hence $\cA$ is standardly stratified, with $\Delta(m)=\Delta_m$; and for $s<m$ the objects $\Delta(s)$ and $\nablabar(s)$ are those of $\cA_{m-1}$, namely $\Delta_s$ and $\nablabar_s$.
\item By \cref{prop:step}(4), $\nablabar_m$ embeds into $I(m)$, and no subobject of $I(m)$ properly containing $\nablabar_m$ has composition factors $L(t)$, $t<m$, besides a single $L(m)$. Indeed, such a subobject would contain a non-split extension of some $L(t)$, $t<m$, by $\nablabar_m$. So $\nablabar(m)=\nablabar_m$.
\end{itemize}
This proves $(\mathrm I_m)$, and hence (1) and (2). The last assertion of (2) is \cref{lem:strat-basic}(2).

\emph{Sufficiency in (3).} Under \eqref{eq:starstar} the same induction with $r=\infty$ gives $\Hom_\cD(G,Y[l])=0$ for all $Y\in\cA$ and $l\ge1$. Let $X,Y\in\cA$, $n\ge2$ and $f\in\Hom(X,Y[n])$. Choose an epimorphism $G^{\oplus k}\twoheadrightarrow X$; its composite with $f$ vanishes. So \cref{thm:realisation} shows that $\real$ is fully faithful, hence an equivalence. Moreover:
\begin{itemize}
\item $G$ satisfies $\Hom(G,G[l])=0$ for $l\ne0$, and by induction (using \eqref{eq:univ-tri}) $\thick G=\thick{\Delta_1,\dots,\Delta_m}$.
\item Since $\real$ is the identity on $\cA$, it identifies $\Kb(\proj\cA)=\thick G_{\Db(\cA)}$ with $\thick{\Delta_1,\dots,\Delta_m}$. This equals $\cD$ if and only if $\gldim\cA<\infty$, i.e.\ (by \cref{thm:C}(4)) if and only if $\cA$ is quasi-hereditary.
\item The equivalence $\cT_s\simeq\Db(\modd B_s)$ is \cref{lem:local-stratum-derived}.
\end{itemize}

\emph{Necessity in (3).} If $\real$ is an equivalence then, by \cref{lem:hom-duality}, (2) and \cref{lem:strat-basic}(3),
\[\Hom^l_{\cT_s}(\Pi_s,S_s)\cong\Hom^l_\cD(\Delta_s,\nablabar_s)\cong\Ext^l_\cA(\Delta(s),\nablabar(s))=0\qquad\text{for }l\ge1.\qedhere\]
\end{proof}

\begin{remark}[The induction in terms of algebras]
Unwinding the induction, the chain $\cD_1\subset\dots\subset\cD_m$ gives a chain of Serre subcategories $\cA_1\subset\dots\subset\cA_m=\cA$. Under $\cA\simeq\modd\Lambda$ this corresponds to a chain of idempotent ideals of $\Lambda$. By \cref{prop:step}(3), each of these ideals is projective as a right module over the corresponding quotient algebra, i.e.\ it is a ``standardly stratifying'' ideal in the sense of \cref{rem:heredity}. The stratum algebras are the $B_s=\End(\Pi_s)$. For $B_s=\kk$ this is exactly a heredity chain.
\end{remark}


\begin{corollary}[cf.\ {\cref{thm:BB-3.4}(2)} and {\cref{thm:BB-3.8}}]\label{cor:forward}
Let $\cA$ be an abelian category. Suppose $\Db(\cA)$ admits an admissible filtration as in the projective case of \cref{set:local} such that:
\begin{itemize}
\item the t-structures on the strata are bounded with local hearts;
\item $\Delta_s,\nablabar_s\in\cA$ for all $s$;
\item \eqref{eq:star} holds.
\end{itemize}
Then $\cA$ is standardly stratified, with standard objects $\Delta_s$ and proper costandard objects $\nablabar_s$. Moreover \eqref{eq:starstar} holds, and $\cT_s\simeq\Db(\modd B_s)$ for all $s$.
\end{corollary}

\begin{proof}
By \cref{prop:recognise}, the standard t-structure of $\Db(\cA)$ is the glued one, so \cref{thm:A}(1),(2) apply. The identity functor of $\Db(\cA)$ is a realisation functor for the standard t-structure, so \cref{thm:A}(3) gives the remaining statements.
\end{proof}

\subsection{Comparison with known results}

\begin{remark}[Bodzenta--Bondal]\label{rem:compare-BB}
In the exceptional case (\cref{ex:local-settings}(1)) we have $\Pi_s=S_s=\kk$, $B_s=\kk$, and \eqref{eq:starstar} holds trivially. Then:
\begin{itemize}
\item \cref{prop:criterion} together with \cref{thm:A}(1),(2) is \cref{thm:BB-3.1}(1). Indeed, standardly stratified with $B_s=\kk$ means quasi-hereditary, i.e.\ highest weight, and for quasi-hereditary algebras $\nablabar=\nabla$.
\item \cref{thm:A}(3) is \cref{thm:BB-3.1}(2), with the hypothesis ``dg enhanced and idempotent complete'' replaced by ``algebraic''.
\item \cref{prop:aisles} is \cref{thm:BB-2.4}, and \cref{prop:recognise} and \cref{cor:forward} are the corresponding parts of \cref{thm:BB-3.4,thm:BB-3.8}.
\end{itemize}
The proof of \cref{thm:A} differs from that of \cite{BB26}. Bodzenta and Bondal first show that $G$ is a tilting generator of $\cD$, and then compare t-structures. When $B_s\neq\kk$, the object $G$ generates only $\thick{\Delta_1,\dots,\Delta_m}\subsetneq\cD$ (\cref{thm:C}), so this route is unavailable.

Instead we proceed by induction on the number of strata, as in the theory of heredity chains. \Cref{prop:step} establishes Ext-projectivity in degrees $1$ and $2$ directly, and \cref{lem:strat-induction} assembles the steps. This is why the degree-two condition \eqref{eq:star} appears.
\end{remark}

\begin{remark}[Stratifying systems]
The sequence $(\Delta_1,\dots,\Delta_m)$ in $\cA$ satisfies the following conditions: $\Hom(\Delta_t,\Delta_s)=0$ for $t>s$, $\Ext^1(\Delta_t,\Delta_s)=0$ for $t\ge s$, and each $\End(\Delta_s)$ is local. These are the axioms of a \emph{stratifying system} in the sense of Erdmann and S\'aenz \cite{ES03}; see also \cite{MMS04}. The theory of stratifying systems generalises the Dlab--Ringel standardisation theorem (\cref{thm:DR}) from quasi-hereditary to standardly stratified algebras, by passing to the endomorphism algebra of a relative projective generator. \Cref{thm:A}(1),(2) says that in the glued situation no such passage is needed, in the same way that \cref{thm:BB-3.1} sharpens \cref{thm:DR}.
\end{remark}

\begin{remark}[On the role of \eqref{eq:star}]\label{rem:star}
In \cref{prop:step}, condition (ii) with $l=2$ enters only through step (c), via $V_s\otimes\Hom_\cT(\Pi,S[2])$. For the top stratum it can therefore be dropped if $V_s=\Ext^1_\cA(G'_s,\Delta)=0$ for all $s$.

Condition (iii) with $l=2$, i.e.\ \eqref{eq:star} for the earlier strata via the induction, enters in two places: in step (d), and hence in the case $l=1$ (step (e)); and in part (3). \Cref{ex:k2} shows that the conclusion of \cref{thm:A}(1) may hold even when \eqref{eq:star} fails.
\end{remark}

% ===============================================================
\section{A stratified Bodzenta--Bondal criterion}\label{sec:criterion}
% ===============================================================

\subsection{Hom-dual triples}

\begin{definition}\label{def:hom-dual}
Let $(E_1,\dots,E_m)$ be a local-exceptional sequence in $\cD$.
\begin{itemize}
\item A \emph{Hom-dual sequence} for it is a sequence $(N_m,\dots,N_1)$ of objects such that:
\begin{itemize}
\item[\textup{(HD1)}] $\Hom^\bullet(E_s,N_t)\cong\delta_{st}\kk$, concentrated in degree $0$, for all $s,t$;
\item[\textup{(HD2)}] $\Hom^\bullet(N_s,N_t)=0$ for $s<t$; and for all $s$, $\End(N_s)=\kk$ and $\Hom(N_s,N_s[l])=0$ for $l<0$.
\end{itemize}
Put $\cD_s:=\thick{N_1,\dots,N_s}$.
\item A \emph{Hom-dual triple} $(E_\bullet,\Deltabar_\bullet,N_\bullet)$ consists of $(E_1,\dots,E_m)$, a Hom-dual sequence $(N_m,\dots,N_1)$, and objects $\Deltabar_1,\dots,\Deltabar_m$ with morphisms $\delta_s\colon\Deltabar_s\to N_s$, such that:
\begin{itemize}
\item[\textup{(HD3)}] $\cD_m=\cD$;
\item[\textup{(HD4)}] for every $s$, the cone of $\delta_s$ lies in $\cD_{s-1}$, and $\Hom^\bullet(\Deltabar_s,N_t)=0$ for all $t<s$.
\end{itemize}
We call $(\Deltabar_1,\dots,\Deltabar_m)$ a \emph{right dual} of $(N_m,\dots,N_1)$; see \cref{rem:HD4} for the reason.
\end{itemize}
\end{definition}

\begin{intuition}
In the strata of the filtration $\cD_s$ (\cref{lem:peel}), the three sequences play the following roles:
\begin{itemize}
\item $E_s$ is the projective object of the $s$-th stratum, extended by $j_{s!}$;
\item $N_s$ is the simple object, extended by $j_{s*}$;
\item $\Deltabar_s$ is the simple object, extended by $j_{s!}$.
\end{itemize}
There are thus two dualities. $N_\bullet$ and $\Deltabar_\bullet$ are related by mutation, as in Bondal's theory of exceptional sequences. $E_\bullet$ and $N_\bullet$ are related by the Hom-duality (HD1), which pairs a projective with a simple, as in Koszul duality.

In the exceptional case projective and simple objects of the strata coincide, so $\Deltabar_\bullet=E_\bullet$ and the second duality is the first (\cref{prop:exc-dual}). For a standardly stratified algebra one expects the triple (standard, proper standard, proper costandard); see \cref{prop:dualisable}.
\end{intuition}


\begin{lemma}\label{lem:sod}
Let $\cC=\thick{\cU\cup\cV}$ for thick subcategories $\cU,\cV\subset\cC$ with $\Hom^\bullet(\cV,\cU)=0$.
\begin{enumerate}
\item Every $X\in\cC$ fits into a triangle $V_X\to X\to U_X\to V_X[1]$ with $V_X\in\cV$ and $U_X\in\cU$.
\item $X\mapsto V_X$ is right adjoint to the inclusion $\cV\subset\cC$, and $X\mapsto U_X$ is left adjoint to the inclusion $\cU\subset\cC$. Moreover $\cV^\perp\cap\cC=\cU$ and ${}^\perp\cU\cap\cC=\cV$.
\item If $\Hom^\bullet(\cV,X)=0$ then $X\in\cU$, and if $\Hom^\bullet(X,\cU)=0$ then $X\in\cV$.
\end{enumerate}
\end{lemma}

\begin{proof}
Let $\cC'\subset\cC$ be the full subcategory of objects admitting a triangle as in (1).

For $X\in\cC'$ and $V\in\cV$, the vanishing of $\Hom^\bullet(V,U_X)$ gives $\Hom(V,V_X)\cong\Hom(V,X)$. Dually, $\Hom(U_X,U)\cong\Hom(X,U)$ for $U\in\cU$. Hence the triangle is unique up to unique isomorphism, and it is functorial in $X\in\cC'$.

The subcategory $\cC'$ contains $\cU$ and $\cV$ and is closed under shifts. It is closed under cones by the $3\times3$ lemma \cite[Prop.~1.1.11]{BBD82}, applied to the morphism of triangles induced by a morphism in $\cC'$. It is closed under direct summands, because idempotents act on the functorial triangle and split in the thick subcategories $\cU$ and $\cV$. Hence $\cC'=\cC$, which proves (1); the adjunctions in (2) were shown above.

For (3): if $\Hom^\bullet(\cV,X)=0$, then $V_X\to X$ vanishes, so $U_X\cong X\oplus V_X[1]$. Then $V_X[1]$ lies in $\cU\cap\cV$, which is $0$ because $\Hom(\cV,\cU)=0$. Hence $X\cong U_X\in\cU$. The second statement is dual.

The orthogonality statements in (2) follow from (3).
\end{proof}

\begin{lemma}\label{lem:right-dual}
Assume (HD2), (HD3) and (HD4). For every $s$:
\begin{enumerate}
\item $\Deltabar_s\in\cD_s$, $\Hom^\bullet(\Deltabar_s,\cD_{s-1})=0$, and $\cD_s=\thick{\cD_{s-1}\cup\{\Deltabar_s\}}$;
\item $\Hom^\bullet(\Deltabar_s,N_t)=0$ for $t\ne s$, and composition with $\delta_s$ induces an isomorphism $\Hom^\bullet(\Deltabar_s,N_s)\cong\End^\bullet(N_s)$;
\item the pair $(\Deltabar_s,\delta_s)$ is unique up to unique isomorphism.
\end{enumerate}
\end{lemma}

\begin{proof}
(1) The cone of $\delta_s$ and $N_s$ lie in $\cD_s$, hence so does $\Deltabar_s$. Conversely $N_s\in\thick{\cD_{s-1}\cup\{\Deltabar_s\}}$, which gives $\cD_s=\thick{\cD_{s-1}\cup\{\Deltabar_s\}}$. The vanishing $\Hom^\bullet(\Deltabar_s,\cD_{s-1})=0$ follows from (HD4), since $\cD_{s-1}$ is the thick closure of $N_1,\dots,N_{s-1}$.

(2) For $t>s$ the vanishing holds because $\Deltabar_s\in\cD_s$ and $\Hom^\bullet(\cD_s,N_t)=0$ by (HD2). For $t=s$, apply $\Hom^\bullet(-,N_s)$ to the triangle $\Deltabar_s\to N_s\to\cone(\delta_s)\to$, and use $\Hom^\bullet(\cD_{s-1},N_s)=0$.

(3) Let $(\Deltabar'_s,\delta'_s)$ be a second choice. Since $\Hom^\bullet(\Deltabar_s,\cone(\delta'_s))=0$ by (1), composition with $\delta'_s$ gives $\Hom(\Deltabar_s,\Deltabar'_s)\cong\Hom(\Deltabar_s,N_s)$. So $\delta_s=\delta'_s\varphi$ for a unique $\varphi$. Symmetrically $\delta'_s=\delta_s\psi$, and uniqueness shows that $\varphi\psi$ and $\psi\varphi$ are identities.
\end{proof}

\begin{lemma}[Peeling off the last stratum]\label{lem:peel}
Let $(E_\bullet,\Deltabar_\bullet,N_\bullet)$ be a Hom-dual triple in $\cD$, and put $\cD'=\cD_{m-1}$.
\begin{enumerate}
\item $\cD'$ is admissible in $\cD$. More precisely, \cref{lem:sod} applied to $(\cV,\cU)=(\thick{\Deltabar_m},\cD')$ gives the left adjoint of the inclusion, and applied to $(\cV,\cU)=(\cD',\thick{N_m})$ gives the right adjoint. Moreover ${}^\perp\cD'=\thick{\Deltabar_m}$ and $\cD'^\perp=\thick{N_m}$.
\item For $s<m$ we have $E_s,\Deltabar_s\in\cD'$, and $(E_{<m},\Deltabar_{<m},N_{<m})$ is a Hom-dual triple in $\cD'$.
\item The stratum $\cT=\cD/\cD'$, with the t-structure determined by the simple-minded object $S:=j^*N_m$ (\cref{thm:smc}), is a local stratum of projective type. Its projective cover of $S$ is $\Pi:=j^*E_m$, with $\End(\Pi)\cong\End(E_m)$ and $\Hom_\cT(\Pi,S[l])=0$ for $l\neq0$. Moreover
\[j_!\Pi\cong E_m,\qquad j_!S\cong\Deltabar_m,\qquad j_*S\cong N_m.\]
\end{enumerate}
\end{lemma}

\begin{proof}
(1) We have $\cD=\thick{\cD'\cup\{\Deltabar_m\}}$ with $\Hom^\bullet(\Deltabar_m,\cD')=0$ (\cref{lem:right-dual}). We also have $\cD=\thick{\cD'\cup\{N_m\}}$ with $\Hom^\bullet(\cD',N_m)=0$, by (HD2) and (HD3). Now apply \cref{lem:sod}(2) twice.

(2) Let $s<m$. By (HD1), $\Hom^\bullet(E_s,N_m)=0$, so $E_s\in{}^\perp\thick{N_m}=\cD'$ by \cref{lem:sod}(2). Moreover $\Deltabar_s\in\cD_s\subset\cD'$ by \cref{lem:right-dual}(1). Conditions (HD1), (HD2) and (HD4) for indices $<m$ are inherited, and (HD3) holds for $\cD'$ by definition.

(3) We locate the three objects.
\begin{itemize}
\item By (HD1), $\Hom^\bullet(E_m,\cD')=0$. Since ${}^\perp\cD'$ is the essential image of $j_!$, we get $E_m\cong j_!\Pi$.
\item Similarly $\Deltabar_m\cong j_!j^*\Deltabar_m$, and $j^*\Deltabar_m\cong j^*N_m=S$ because $\cone(\delta_m)\in\cD'$.
\item $N_m\in\cD'^\perp$, the essential image of $j_*$, so $N_m\cong j_*S$.
\end{itemize}
Next, $\cT$ is generated by $S$, and $\End^\bullet_\cT(S)\cong\End^\bullet(N_m)$ because $j_*$ is fully faithful. So $S$ is simple-minded by (HD2), and \cref{thm:smc} shows that $\cT$ is a local stratum with simple object $S$. By adjunction and (HD1),
\[\Hom^\bullet_\cT(\Pi,S)\cong\Hom^\bullet(E_m,N_m)\cong\kk,\]
concentrated in degree $0$. Moreover $\End^\bullet(\Pi)\cong\End^\bullet(E_m)$ is concentrated in degree $0$. Now \cref{lem:recognise-proj} applies.
\end{proof}

\begin{remark}[Admissible subcategories and mutations]\label{rem:HD4}
Assume (HD1)--(HD3). Then (HD4) holds if and only if $\cD_{s-1}$ is left admissible in $\cD_s$ for all $s$.
\begin{itemize}
\item One direction is \cref{lem:peel}(1), applied to $\cD_s$.
\item Conversely, if $\cD_{s-1}$ is left admissible in $\cD_s$, the triangle $X'\to N_s\to i^*N_s$ with $X'\in{}^\perp\cD_{s-1}$ provides $\Deltabar_s:=X'$ and $\delta_s$.
\end{itemize}
In this language $\Deltabar_s$ is the right mutation of $N_s$ through $\cD_{s-1}$, which explains the term ``right dual''. Right admissibility, i.e.\ the existence of $j_{s*}$, is automatic by (HD2) and \cref{lem:sod}.

So (HD4) is the only condition of admissibility type, and it asks precisely for the objects $j_{s!}S_s$. On the subcategory $\thick{\Pi_s}\simeq\Kb(\proj B_s)$ of the stratum, the functor $j_{s!}$ is given by $E_s$, but $\Pi_s$ does not generate the stratum unless $B_s=\kk$. A one-sided gluing in the spirit of Collins and Polishchuk \cite{CP10} might allow one to drop (HD4); we have not checked whether their hypotheses apply.
\end{remark}

\begin{proposition}\label{prop:exc-dual}
Let $(E_1,\dots,E_n)$ be a full exceptional sequence in $\cD$ with left dual $(F_n,\dots,F_1)$. Then $(E_\bullet,E_\bullet,F_\bullet)$ is a Hom-dual triple: the left dual is a Hom-dual sequence, and its right dual is the original sequence.
\end{proposition}

\begin{proof}
\textbf{(HD1)} is \eqref{eq:hom-duality-exc}.

\textbf{(HD2)} The left dual is an exceptional sequence \cite{Bon89}.

\textbf{(HD3)} $F_s$ is obtained from $E_s$ by mutation through $\thick{E_1,\dots,E_{s-1}}$, via a triangle $i_*i^!E_s\to E_s\to F_s\to$. So induction on $s$ gives $\thick{F_1,\dots,F_s}=\thick{E_1,\dots,E_s}$, and hence (HD3).

\textbf{(HD4)} The same triangle gives a morphism $\delta_s\colon E_s\to F_s$ whose cone $i_*i^!E_s[1]$ lies in $\cD_{s-1}$. By (HD1), $\Hom^\bullet(E_s,F_t)=0$ for $t<s$.
\end{proof}

\subsection{The criterion}

\begin{theorem}\label{thm:B}
Let $\cD$ be a Hom-finite idempotent complete triangulated category, and let $(E_\bullet,\Deltabar_\bullet,N_\bullet)$ be a Hom-dual triple in $\cD$. Put $B_s:=\End(E_s)$.
\begin{enumerate}
\item $\cD_s=\thick{N_1,\dots,N_s}$ is an admissible filtration of $\cD$. Each stratum $\cT_s$, with the t-structure determined by the simple-minded object $S_s:=j_s^*N_s$, is a local stratum of projective type, and:
\begin{itemize}
\item its projective object is $\Pi_s:=j_s^*E_s$, with $\End(\Pi_s)\cong B_s$ and $\Hom(\Pi_s,S_s[l])=0$ for $l\ne0$;
\item $E_s\cong j_{s!}\Pi_s$, $\Deltabar_s\cong j_{s!}S_s$ and $N_s\cong j_{s*}S_s$.
\end{itemize}
\item Let $\cA$ be the heart of the glued t-structure. Then $\cA$ is standardly stratified with respect to $1<\dots<m$, with standard objects $E_s$ and proper costandard objects $N_s$, if and only if
\[\Hom(E_s,E_t[l])=0\quad\text{and}\quad\Hom(N_s,N_t[l])=0\qquad\text{for all $s,t$ and all $l<0$.}\]
In this case $\End(\Delta(s))\cong B_s$, and $\cA$ is quasi-hereditary if and only if all $E_s$ are exceptional.
\item Suppose in addition that $\cD$ is algebraic and the conditions in \textup{(2)} hold. Then:
\begin{itemize}
\item $\real\colon\Db(\cA)\to\cD$ is an equivalence;
\item $\cT_s\simeq\Db(\modd B_s)$ for every $s$;
\item $\thick{E_1,\dots,E_m}\simeq\Kb(\proj\cA)$, and this equals $\cD$ if and only if all $E_s$ are exceptional.
\end{itemize}
\end{enumerate}
\end{theorem}

\begin{proof}
Everything is proved by induction on $m$, peeling off the last stratum with \cref{lem:peel}.

\emph{Proof of (1).} By \cref{lem:peel}(2), applied $m-s$ times, the truncated triple $(E_{\le s},\Deltabar_{\le s},N_{\le s})$ is a Hom-dual triple in $\cD_s$. Applying \cref{lem:peel}(1),(3) to it gives all claims about $\cD_{s-1}\subset\cD_s$ and $\cT_s$.

\emph{Setup for (2) and (3).} By (1) we are in the projective case of \cref{set:local}, with
\[\Delta_s=E_s,\qquad\nablabar_s=N_s,\]
and \eqref{eq:starstar} holds for every stratum. By \cref{prop:criterion}, the vanishing conditions in (2) are equivalent to $E_s,N_s\in\cA$ for all $s$; conversely, if $\cA$ is standardly stratified with these objects, the conditions hold. So assume $E_s,N_s\in\cA$ for all $s$, and prove by induction on $m$:
\begin{quote}
$(\mathrm{I}_m)$: $\cA$ is standardly stratified with $\Delta(s)=E_s$ and $\nablabar(s)=N_s$, and the iterated universal extension $G$ of $(E_1,\dots,E_m)$ is a projective generator of $\cA$ with $\Hom(G,Y[l])=0$ for all $Y\in\cA$ and $l\ge1$.
\end{quote}

\emph{Base case $m=1$.} By \cref{lem:peel}(3), $\cD=\cT_1$ is a local stratum of projective type with $\Pi=E_1$ and $S=N_1$, satisfying \eqref{eq:starstar}. So $\cA\simeq\modd B_1$, and the vanishing is \cref{lem:local-stratum-derived}(1).

\emph{Inductive step $m\ge2$.} The truncated triple in $\cD'=\cD_{m-1}$ is a Hom-dual triple (\cref{lem:peel}(2)). Its glued heart is $\cA\cap\cD'$, which contains $E_s,N_s$ for $s<m$. So $(\mathrm I_{m-1})$ holds for it. By \cref{lem:peel}(3), the recollement $(\cD',\cD,\cT_m)$ satisfies the hypotheses of \cref{prop:step} with $r=\infty$, $\Delta=E_m$ and $\nablabar=N_m$. Now argue exactly as in the inductive step of the proof of \cref{thm:A}: \cref{prop:step} gives the projective generator and the trace condition, and \cref{lem:strat-induction} gives the standard stratification. This proves $(\mathrm I_m)$, hence (2); the last assertion of (2) is \cref{lem:strat-basic}(2).

\emph{Proof of (3).} By $(\mathrm I_m)$ and \cref{thm:realisation}, $\real$ is an equivalence, as in the proof of \cref{thm:A}(3). The remaining statements are \cref{thm:A}(3) and \cref{thm:C}(4).
\end{proof}

\begin{remark}[Comparison with Bodzenta--Bondal]\label{rem:B-vs-BB}
For a full exceptional sequence and the triple $(E_\bullet,E_\bullet,F_\bullet)$ of \cref{prop:exc-dual}, all strata are $\Db(\kk)$. Then \cref{thm:B}(2),(3) become \cref{thm:BB-3.1}(1),(2), with ``dg enhanced'' replaced by ``algebraic''. The new features in the general case are the following.
\begin{enumerate}
\item The Hom-dual sequence is extra data: the left dual pairs with $E_s$ via $\End(E_s)$ rather than $\kk$.
\item The ambient category is generated by the Hom-dual sequence; the $E_s$ generate only $\Kb(\proj\cA)$ (\cref{thm:C}).
\item The right dual of $N_\bullet$ is a third sequence $\Deltabar_\bullet$, of ``proper standard'' objects, instead of $E_\bullet$ itself.
\end{enumerate}
The structure of the argument is the same as for quasi-hereditary algebras. \Cref{lem:peel} splits off the last stratum, \cref{prop:step} shows that the corresponding trace ideal is projective, and \cref{lem:strat-induction} assembles the steps.
\end{remark}


% ===============================================================
\section{Standardly stratified algebras}\label{sec:stratified}
% ===============================================================

\subsection{Standard modules form local-exceptional sequences}

\begin{theorem}\label{thm:C}
Let $\Lambda$ be a standardly stratified algebra with respect to $1<\dots<m$, and let $B_i=\End(\Delta(i))$.
\begin{enumerate}
\item $(\Delta(1),\dots,\Delta(m))$ is a full local-exceptional sequence in $\Kb(\proj\Lambda)$ with $\End(\Delta(i))=B_i$.
\item In $\Db(\modd\Lambda)$ we have $\Hom^\bullet(\Delta(i),\nablabar(i))\cong\kk$ (in degree $0$), and $\Hom^\bullet(\Delta(j),\nablabar(i))=0$ for $j>i$.
\item If $B_i\ne\kk$, then $\nablabar(i)\notin\Kb(\proj\Lambda)$, i.e.\ $\pd\nablabar(i)=\infty$.
\item $\gldim\Lambda<\infty$ if and only if $\Lambda$ is quasi-hereditary, if and only if $B_i\cong\kk$ for all $i$.
\end{enumerate}
\end{theorem}

\begin{proof}
(1) and (2) are \cref{lem:strat-basic}.

(3) Suppose $X:=\nablabar(i)\in\Kb(\proj\Lambda)=\thick{\Delta(1),\dots,\Delta(m)}$.
\begin{itemize}
\item Apply \cref{lem:sod}(3) with $\cU=\thick{\Delta(1),\dots,\Delta(i)}$ and $\cV=\thick{\Delta(i+1),\dots,\Delta(m)}$. Since $\Hom^\bullet(\cV,X)=0$ by~(2), we get $X\in\cU$.
\item Apply \cref{lem:sod}(1) to $\cU$, with $\cU'=\thick{\Delta(1),\dots,\Delta(i-1)}$ and $\cV'=\thick{\Delta(i)}$. This gives a triangle $V\to X\to U\to V[1]$ with $V\in\thick{\Delta(i)}$ and $U\in\cU'$.
\item Since $\Hom^\bullet(\Delta(i),\cU')=0$, we get $\Hom^\bullet(\Delta(i),V)\cong\Hom^\bullet(\Delta(i),X)\cong\kk$.
\end{itemize}
This contradicts \cref{prop:local-exc}(3) applied to $E=\Delta(i)$ in the algebraic category $\Db(\modd\Lambda)$.

(4) If $\gldim\Lambda<\infty$, every module is perfect, so $B_i=\kk$ for all $i$ by~(3). Then $\Lambda$ is quasi-hereditary by \cref{lem:strat-basic}(2). Conversely, quasi-hereditary algebras have finite global dimension \cite{CPS88,DR92}.
\end{proof}

\begin{remark}
Statement~(4) is presumably well known to experts (cf.\ \cite{AHLU00,Dla96}). We include the short argument because it isolates exactly the obstruction discussed in the introduction. In a genuinely stratified situation, the standard objects generate only $\Kb(\proj\Lambda)$, while their Hom-dual partners $\nablabar(i)$ lie outside it.
\end{remark}

\subsection{The converse}
We now ask which standardly stratified algebras arise from \cref{thm:B}.

Let $\cA_s\subset\modd\Lambda$ be the Serre subcategory of modules whose composition factors lie in $\{L(1),\dots,L(s)\}$. We say that $\Lambda$ satisfies the \emph{$\nablabar$-condition} if, for every $i$, there is an exact sequence
\[0\to\nablabar(i)\to I(i)\to C_i\to0\qquad\text{with } C_i\in\cF(\nablabar(t):t>i).\]

\begin{lemma}\label{lem:nablabar}
Let $\Lambda$ be standardly stratified.
\begin{enumerate}
\item $\End(\nablabar(i))=\kk$.
\item $\thick{\nablabar(1),\dots,\nablabar(s)}=\Db_{\cA_s}(\modd\Lambda)$, the subcategory of complexes with cohomology in $\cA_s$. In particular $\Db(\modd\Lambda)=\thick{\nablabar(1),\dots,\nablabar(m)}$.
\item If $\Lambda$ satisfies the $\nablabar$-condition, then $(\nablabar(m),\dots,\nablabar(1))$ is a Hom-dual sequence for $(\Delta(1),\dots,\Delta(m))$ in $\Db(\modd\Lambda)$ satisfying (HD3); i.e.\ (HD1)--(HD3) hold.
\end{enumerate}
\end{lemma}

\begin{proof}
(1) Let $f$ be a non-zero endomorphism of $\nablabar(i)$. If $\ker f$ contained $\soc\nablabar(i)=L(i)$, then $\im f\ne0$ would have no composition factor $L(i)$, since $[\nablabar(i):L(i)]=1$. But $\im f$ contains the socle $L(i)$, a contradiction. Hence $\ker f=0$ (the socle is essential), so $f$ is an automorphism. Thus $\End(\nablabar(i))$ is a division algebra, hence equal to $\kk$.

(2) The subcategory $\Db_{\cA_s}(\modd\Lambda)$ is $\thick{L(1),\dots,L(s)}$. Moreover $\nablabar(t)$ has composition factors $L(t)$ (once) and $L(u)$ with $u<t$. By induction on $s$, $\thick{\nablabar(1),\dots,\nablabar(s)}=\thick{L(1),\dots,L(s)}$.

(3) We first prove the dual of \cref{lem:strat-basic}(1): if $X\in\cF(\nablabar(t):t>i)$ and $Y$ has composition factors in $\{L(u):u\le i\}$, then $\Ext^\bullet(Y,X)=0$. It suffices to take $X=\nablabar(t)$ with $t>i$, and we use descending induction on $t$.
\begin{itemize}
\item $\Hom(Y,\nablabar(t))=0$, since the socle $L(t)$ is not a composition factor of $Y$.
\item $\Hom(Y,I(t))=0$ for the same reason.
\item The sequence $0\to\nablabar(t)\to I(t)\to C_t\to0$ gives $\Ext^1(Y,\nablabar(t))\cong\Hom(Y,C_t)$ and $\Ext^{l+1}(Y,\nablabar(t))\cong\Ext^l(Y,C_t)$ for $l\ge1$.
\item These groups vanish by induction, since $C_t\in\cF(\nablabar(u):u>t)$.
\end{itemize}
Now we check the conditions.
\begin{itemize}
\item (HD1) for $j\ge i$ is \cref{lem:strat-basic}(3). For $j<i$, we have $\Hom(\Delta(j),\nablabar(i))=0$ by \cref{lem:strat-basic}(3). The sequence $0\to\nablabar(i)\to I(i)\to C_i\to0$ gives $\Ext^1(\Delta(j),\nablabar(i))\cong\Hom(\Delta(j),C_i)$ and $\Ext^{l+1}(\Delta(j),\nablabar(i))\cong\Ext^l(\Delta(j),C_i)$ for $l\ge1$. These vanish by the dual statement, because $\Delta(j)$ has composition factors $L(u)$ with $u\le j<i$.
\item (HD2): for $s<t$, take $X=\nablabar(t)$ and $Y=\nablabar(s)$ in the dual statement. The remaining parts of (HD2) are (1) and the fact that modules have no negative extensions.
\item (HD3) is (2). \qedhere
\end{itemize}
\end{proof}

\begin{remark}
The $\nablabar$-condition is the injective counterpart of standard stratification. We expect that it holds for every standardly stratified algebra: it is a form of BGG reciprocity, cf.\ \cite{AHLU00,Dla96}. We have not checked this and do not use it. For quasi-hereditary algebras it is the classical fact that injective modules have $\nabla$-filtrations \cite{CPS88,DR92}.
\end{remark}

\begin{proposition}\label{prop:dualisable}
Let $\Lambda$ be standardly stratified and satisfy the $\nablabar$-condition, and suppose that
\[\Ext^l_\Lambda(\Deltabar(s),\nablabar(t))=0\qquad\text{for all }t<s\text{ and all }l.\]
Let $\delta_s$ be the composite $\Deltabar(s)\twoheadrightarrow L(s)\hookrightarrow\nablabar(s)$. Then $\big((\Delta(s))_s,(\Deltabar(s))_s,(\nablabar(s))_s\big)$, with the morphisms $\delta_s$, is a Hom-dual triple in $\Db(\modd\Lambda)$.
\end{proposition}

\begin{proof}
By \cref{lem:nablabar}(3), conditions (HD1)--(HD3) hold. For (HD4), recall that $[\Deltabar(s):L(s)]=1=[\nablabar(s):L(s)]$, and that all other composition factors of $\Deltabar(s)$ and $\nablabar(s)$ are $L(t)$ with $t<s$. So the kernel $\rad\Deltabar(s)$ and the cokernel $\nablabar(s)/L(s)$ of $\delta_s$ lie in $\cA_{s-1}$. Hence $\cone(\delta_s)$, whose cohomology consists of these two modules, lies in $\Db_{\cA_{s-1}}(\modd\Lambda)=\cD_{s-1}$, by \cref{lem:nablabar}(2). The Hom condition in (HD4) is the hypothesis.
\end{proof}

\begin{definition}
Let $\cA$ be a $\kk$-linear abelian category such that $\Db(\cA)$ is Hom-finite.
\begin{itemize}
\item A \emph{standardly stratified structure} on $\cA$ is a total order $L_1<\dots<L_m$ on the isomorphism classes of simple objects for which $\cA$ is Deligne finite and standardly stratified.
\item It is \emph{dualisable} if there are objects $\Deltabar_s\in\Db(\cA)$ and morphisms $\delta_s$ making $\big((\Delta(s))_s,(\Deltabar_s)_s,(\nablabar(s))_s\big)$ a Hom-dual triple in $\Db(\cA)$. By \cref{lem:right-dual}(3), the $\Deltabar_s$ are then unique.
\end{itemize}
\end{definition}

\begin{corollary}[Analogue of {\cref{thm:BB-3.8}}]\label{cor:BB38}
Let $\cA$ be a $\kk$-linear abelian category such that $\Db(\cA)$ is Hom-finite. The assignment
\[(L_1<\dots<L_m)\longmapsto\big((\Delta(1),\dots,\Delta(m)),\,(\nablabar(m),\dots,\nablabar(1))\big)\]
is a bijection between:
\begin{enumerate}
\item dualisable standardly stratified structures on $\cA$; and
\item isomorphism classes of pairs $(E_\bullet,N_\bullet)$ of objects of $\cA$ that extend to a Hom-dual triple $(E_\bullet,\Deltabar_\bullet,N_\bullet)$ in $\Db(\cA)$.
\end{enumerate}
In particular every dualisable standardly stratified structure is recovered by gluing. For exceptional sequences this is \cref{thm:BB-3.8}.
\end{corollary}

\begin{proof}
By \cref{thm:C}(1) and the definition, the assignment maps (1) to (2). It is injective, since the order is recovered from the sequence: $L_s$ is the top of $\Delta(s)$.

For surjectivity, let $(E_\bullet,\Deltabar_\bullet,N_\bullet)$ be a Hom-dual triple with $E_s,N_s\in\cA$.
\begin{itemize}
\item By \cref{thm:B}(1) we are in \cref{set:local}, with $\Delta_s=E_s$ and $\nablabar_s=N_s$.
\item Since $E_s,N_s\in\cA$, we have $E_s\in\Db(\cA)^{\le0}$ and $N_s\in\Db(\cA)^{\ge0}$. So, by \cref{prop:recognise}, the standard t-structure is the glued one.
\item Objects of $\cA$ have no negative extensions. So \cref{thm:B}(2) shows that $\cA$ is standardly stratified for the order $L_1<\dots<L_m$, where $L_s$ is the top of $E_s$, with $\Delta(s)=E_s$ and $\nablabar(s)=N_s$.
\end{itemize}
This structure is dualisable via $\Deltabar_\bullet$, and it maps to $(E_\bullet,N_\bullet)$.

For a quasi-hereditary structure, the standard objects form a full exceptional sequence in $\Db(\cA)$ whose left dual consists of the costandard objects $\nabla(s)=\nablabar(s)$ \cite[\S3]{BB26}. So the structure is dualisable with $\Deltabar_\bullet=\Delta_\bullet$, by \cref{prop:exc-dual}, and we recover \cref{thm:BB-3.8}.
\end{proof}

\begin{question}\label{q:HD4}
Is every standardly stratified structure dualisable? More concretely:
\begin{enumerate}
\item[(a)] Does every standardly stratified algebra satisfy the $\nablabar$-condition?
\item[(b)] Is $\Ext^\bullet_\Lambda(\Deltabar(s),\nablabar(t))=0$ for all $t<s$?
\end{enumerate}
By \cref{prop:dualisable}, positive answers to (a) and (b) imply dualisability, with the proper standard modules as right dual.
\end{question}

Some evidence and remarks:
\begin{itemize}
\item For quasi-hereditary algebras, $\Deltabar=\Delta$, $\nablabar=\nabla$, and both (a) and (b) are classical.
\item For properly stratified algebras we expect (b) to hold. On the top stratum, $\Deltabar(m)$ should be $j_{m!}$ of the simple $B_m$-module, and then $\Ext^\bullet(\Deltabar(m),i_*Y)=0$ by adjunction; one would then proceed by induction, peeling off the top stratum.
\item If (b) fails, the structure may still be dualisable, with $\Deltabar_s$ a complex rather than a module.
\end{itemize}

% ===============================================================
\section{Application: \texorpdfstring{$\mathbb{P}^n[k]$}{Pn[k]}-like objects}\label{sec:P-objects}
% ===============================================================

$\Pn nk$-like objects generalise exceptional objects by allowing self-extensions in the degrees $k,2k,\dots,nk$. In this section we explain how they fit into the general framework.
\begin{itemize}
\item For $k=0$ they are local-exceptional objects (\cref{sec:P-zero}), so \cref{thm:B,thm:C} apply.
\item For $k\ge1$ they are simple-minded and generate local strata (\cref{sec:P-positive}), so \cref{thm:A} applies to filtrations by sequences of them (\cref{sec:P-sequences}).
\end{itemize}
\Cref{tab:summary} summarises the outcome.

\begin{table}[ht]
\centering
\renewcommand{\arraystretch}{1.25}
\small
\begin{tabular}{@{}llll@{}}
\toprule
stratum & heart of the stratum & glued heart $\cA$ & $\Db(\cA)\simeq\cD$? \\
\midrule
$\Db(\kk)$ (exceptional) & $\modd\kk$ & highest weight \cite{BB26} & yes \\
$\Db(\modd B)$, $B$ local & $\modd B$ & std.\ stratified (Thm.~\ref{thm:A}, \ref{thm:B}) & yes \\
$\thick P$, $\Pn nk$, $k\ge 3$ & $\modd\kk$ & highest weight (Thm.~\ref{thm:Pseq}) & no \\
$\thick P$, $\Pn nk$, $k=2$ & $\modd\kk$ & open (Question~\ref{q:k2}) & no$^{\ast}$ \\
$\thick P$, $\Pn n1$, $n\ge 3$ & $\modd\kk[x]/(x^2)$ & std.\ stratified (Thm.~\ref{thm:Pseq}) & no \\
$\thick P$, $\Pn 21$ & $\modd\kk[x]/(x^2)$ & open (Question~\ref{q:k2}) & no$^{\ast}$ \\
$\thick P$, $\Pn 11$ & nilpotent $\kk[x]$-modules & not Deligne finite (Thm.~\ref{thm:Pseq}) & --- \\
$\thick P$, $\Pn n0$ & no t-structure \cite{Bonf26} & --- (use Thm.~\ref{thm:B} instead) & --- \\
\bottomrule
\end{tabular}
\caption{Strata of an admissible filtration and the resulting glued heart. In the rows with $\Pn nk$ we assume $n\ge1$. ``Std.\ stratified'' means standardly stratified; ``no t-structure'' means no non-trivial t-structure. A $\Pn n0$-like object enters \cref{thm:B} as a local-exceptional object $E_s$, and its stratum is the one generated by the Hom-dual object $N_s$, equivalent to $\Db(\modd\kk[t]/(t^{n+1}))$. The asterisk means ``no, whenever $\cA$ is a module category of the expected type''.}\label{tab:summary}
\end{table}


\subsection{Definitions and examples}

\begin{definition}[{\cite{Kru18}, \cite[Def.~2.1]{HK19}, \cite[Def.~2.1]{Bonf26}}]\label{def:Pnk}
Let $n\ge0$ and $k\in\mathbb{Z}$. An object $P$ of $\cD$ is \emph{$\Pn nk$-like} if there is an isomorphism of graded algebras
\[\End^\bullet(P)\cong\kk[t]/(t^{n+1}),\qquad\deg t=k.\]
It is a \emph{$\Pn nk$-object} if moreover $\Hom^\bullet(P,X)$ is finite-dimensional for all $X$ and $P$ is $nk$-Calabi--Yau, i.e.\ there are isomorphisms $\Hom(P,X)\cong\Hom(X,P[nk])^\vee$, functorial in $X\in\cD$.
\end{definition}

As in \cite{Bonf26} we allow arbitrary $k\in\mathbb{Z}$. For $k=0$ the condition says that $\End(P)\cong\kk[t]/(t^{n+1})$ and $\Hom(P,P[l])=0$ for $l\ne0$.

\begin{example}\label{ex:Pnk}
\leavevmode
\begin{enumerate}
\item The $\Pn0k$-like objects are exactly the exceptional objects.
\item $\Pn1k$-like objects are the \emph{$k$-spherelike} objects. $\Pn1k$-objects are the $k$-spherical objects of Seidel and Thomas \cite{ST01}; in the terminology of \cite{HJY13} they are \emph{$w$-spherical} with $w=k$.
\item $\Pn n2$-objects are the $\mathbb{P}^n$-objects of Huybrechts and Thomas \cite{HT06}. For example, the structure sheaf of a compact hyperk\"ahler manifold of dimension $2n$ is one.
\item Krug \cite{Kru18} introduced the notation $\Pn nk$ and gave examples of varieties whose structure sheaves are $\Pn nk$-objects; see also \cite{HK19}.
\item ($k=0$) If $B=\kk[t]/(t^{n+1})$, then $B$ is a $\Pn n0$-like object of $\Db(\modd B)$. More generally, a module $M$ with $\End(M)\cong\kk[t]/(t^{n+1})$ and $\Ext^{>0}(M,M)=0$ is $\Pn n0$-like. Standard modules of stratified algebras provide examples (\cref{thm:C}).
\end{enumerate}
\end{example}

\begin{remark}[{\cite[Rem.~2.2]{HK19}, \cite[Lemma~2.5]{Bonf26}}]
A $\Pn nk$-like object $P$ is already a $\Pn nk$-object of the subcategory $\thick P$; in fact, for $\cD$ algebraic, $\thick P$ is $nk$-Calabi--Yau. Hence for questions concerning $\thick P$ alone, the Calabi--Yau condition is irrelevant. In sequences of such objects the Calabi--Yau condition refers to the ambient category and is not preserved by mutations, so we work with ``$\Pn nk$-like'' throughout.
\end{remark}

\begin{intuition}
For an exceptional object $E$ one has $\thick E\simeq\Db(\kk)$, and $E$ is at the same time the simple object, the projective object and the injective object of the heart $\modd\kk$. A $\Pn nk$-like object is a ``thickening'' with self-extensions in degrees $k,2k,\dots,nk$, and these three roles separate:
\begin{itemize}
\item For $k\ge1$, $P$ has no negative self-extensions and $\End(P)=\kk$. It is \emph{simple-minded} and determines a t-structure.
\item For $k\le0$, $P$ has no positive self-extensions, so it is \emph{silting} and determines a co-t-structure. For $k=0$ it behaves like the free module over the local algebra $\kk[t]/(t^{n+1})$.
\end{itemize}
This is the dichotomy of \cref{thm:HJY,thm:Bonfert}; the standard t-structure and the standard weight structure on $\thick P$ are obtained in exactly this way in \cite[\S3.2]{Bonf26}.
\end{intuition}

\subsection{Known classification results}

\begin{theorem}[{Holm--J\o rgensen--Yang \cite{HJY13}}]\label{thm:HJY}
Let $w\in\mathbb{Z}$, and let $\cT$ be the $\kk$-linear algebraic triangulated category generated by a $w$-spherical object.
\begin{itemize}
\item For $w\le0$, $\cT$ has no non-trivial t-structures, but it has one family of non-trivial co-t-structures.
\item For $w\ge1$, $\cT$ has no non-trivial co-t-structures, but it has one family of non-trivial t-structures.
\end{itemize}
For $w=0$, $\cT$ is the compact derived category of the dual numbers.
\end{theorem}

Here ``non-trivial'' excludes $(\cT,0)$ and $(0,\cT)$, and a ``family'' consists of the shifts of a single structure.

\begin{theorem}[{Hochenegger--Krug \cite[Lemma~2.8]{HK19}; Bonfert \cite[Props.~2.3 and~2.4]{Bonf26}}]\label{thm:formality}
For all $n\ge0$ and $k\in\mathbb{Z}$, the graded algebra $A_{n,k}:=\kk[x]/(x^{n+1})$ with $\deg x=k$ is intrinsically formal. Consequently, let $\cD$ be an idempotent complete algebraic triangulated category and $P\in\cD$ a $\Pn nk$-like object. Then $\thick P\simeq\operatorname{perf}(A_{n,k})$, where $A_{n,k}$ is regarded as a dg algebra with zero differential, and $P$ corresponds to $A_{n,k}$.
\end{theorem}

\begin{theorem}[{Bonfert \cite[Thm.~3.3]{Bonf26}}]\label{thm:Bonfert}
Let $\cT$ be the triangulated category generated by a $\Pn nk$-object $P$.
\begin{enumerate}
\item If $n=0$, every non-trivial t-structure on $\cT$ is a shift of the standard t-structure, and every non-trivial weight structure on $\cT$ is a shift of the standard weight structure.
\item If $n>0$ and $k>0$, every non-trivial t-structure on $\cT$ is a shift of the standard t-structure, whose heart is $\Filt(P)$. There is no non-trivial weight structure on $\cT$.
\item If $n>0$ and $k\le0$, every non-trivial weight structure on $\cT$ is a shift of the standard weight structure, whose co-heart is $\add P$. There is no non-trivial t-structure on $\cT$.
\end{enumerate}
\end{theorem}

For $n=1$ this recovers \cref{thm:HJY}. By \cref{thm:formality} and the preceding remark, \cref{thm:Bonfert} applies to $\thick P$ for every $\Pn nk$-like object $P$ of an idempotent complete algebraic triangulated category.

\begin{remark}[{Koszul dual models \cite[\S4]{Bonf26}}]\label{rem:Bonfert-Koszul}
Bonfert also computes the dg Koszul dual $A_{n,k}^{!,\mathrm{dg}}=\RHom_{A_{n,k}}(\kk,\kk)$ \cite[Thm.~4.9, Props.~4.11 and~4.12]{Bonf26}:
\begin{itemize}
\item for $n=1$ and $k\ne1$ it is $\kk[s]$ with $\deg s=1-k$;
\item for $n=k=1$ it is $\kk[[s]]$ with $\deg s=0$;
\item for $n>1$ it is a non-formal dg algebra with cohomology $\kk[s,t]/(s^2)$, where $\deg s=1-k$ and $\deg t=2-(n+1)k$.
\end{itemize}
For $k=0$ this is the Ext-algebra $\Ext^\bullet_B(\kk,\kk)$ of the simple module over $B=\kk[t]/(t^{n+1})$. It is the graded endomorphism algebra of the ``simple-like'' objects $\nablabar$ that will be paired with $\Pn n0$-like objects in \cref{sec:gluing,sec:stratified}.
\end{remark}

\subsection{The case \texorpdfstring{$k=0$}{k=0}: \texorpdfstring{$\mathbb{P}^n[0]$}{Pn[0]}-like objects are local-exceptional}\label{sec:P-zero}

A $\Pn n0$-like object is the same as a local-exceptional object $E$ with $\End(E)\cong\kk[t]/(t^{n+1})$. In this case the general results specialise as follows.
\begin{itemize}
\item \emph{The category generated by $E$.} By \cref{thm:formality} and \cref{thm:Bonfert}(3), $\thick E\simeq\operatorname{perf}(\kk[t]/(t^{n+1}))$ carries, up to shift, a unique non-trivial weight structure, with co-heart $\add E$, and no non-trivial t-structure. This sharpens \cref{prop:local-exc} for this particular local algebra.
\item \emph{Gluing.} Consider a semi-orthogonal sequence of $\Pn{n_s}0$-like objects $(E_1,\dots,E_m)$ which is part of a Hom-dual triple. By \cref{thm:B}, the glued heart is standardly stratified with $B_s\cong\kk[t]/(t^{n_s+1})$ if and only if neither sequence has negative extensions.
\item \emph{Stratified algebras.} Conversely, by \cref{thm:C}, a standard module $\Delta(i)$ of a standardly stratified algebra with $B_i\cong\kk[t]/(t^{n+1})$ is $\Pn n0$-like.
\end{itemize}

The Hom-dual partners are of a different nature. Within its stratum, $N_s$ corresponds to the simple module over $B_s=\kk[t]/(t^{n_s+1})$. Its graded endomorphism algebra is $\Ext^\bullet_{B_s}(\kk,\kk)$, which by \cref{rem:Bonfert-Koszul} (with $k=0$) has the following form:
\begin{itemize}
\item for $n_s=1$, it is $\kk[s]$ with $\deg s=1$, so $N_s$ is ``$\Pn\infty1$-like'';
\item for $n_s\ge2$, its cohomology is $\kk[s,t]/(s^2)$ with $\deg s=1$ and $\deg t=2$, so $N_s$ is not $\mathbb{P}$-like.
\end{itemize}
The simplest example is $m=1$: $E=B$ and $\Deltabar=N=\kk$ in $\Db(\modd B)$, and \cref{thm:B} recovers $\modd B$.


\subsection{The case \texorpdfstring{$k\ge1$}{k>=1}: local strata generated by \texorpdfstring{$\mathbb{P}$}{P}-like objects}\label{sec:P-positive}

For $k\ge1$, a $\Pn nk$-like object has no negative self-extensions and satisfies $\End(P)=\kk$. It is therefore simple-minded, and it generates a local stratum (\cref{def:local-stratum}).


\begin{proposition}[{cf.\ \cite[\S3.2]{Bonf26}}]\label{prop:tP}
Let $P$ be $\Pn nk$-like with $k\ge1$. Then $\{P\}$ is a simple-minded collection in $\thick P$. Consequently $\thick P$ carries a bounded t-structure $t_P$ with heart $\cH_P:=\Filt(P)$, a length category with unique simple object $P$. The aisles of $t_P$ are described by \cref{lem:aisles} with $Q=S=P$.
\end{proposition}

\begin{proof}
Condition (SM1) holds because $\End^\bullet(P)$ is concentrated in degrees $0,k,\dots,nk\ge0$. Condition (SM2) holds because $\End(P)=\kk$ (as $k\ne0$), and (SM3) is clear. Now apply \cref{thm:smc}, and \cref{lem:aisles} with $Q=S=P$.
\end{proof}

\begin{remark}\label{rem:tP-unique}
The t-structure $t_P$ is the standard t-structure of the proper positive dg algebra $A_{n,k}$ \cite[Thm.~7.1]{KN13}, \cite[\S3.2]{Bonf26}. By \cref{thm:Bonfert}(2), every non-trivial t-structure on $\thick P$ is a shift of $t_P$. Hence, for strata generated by $\Pn nk$-like objects with $k\ge1$, the gluing data used in \cref{sec:gluing,sec:P-sequences} are determined up to shifts.
\end{remark}

\begin{proposition}[{cf.\ \cite[Lemmas~3.11 and~4.15, Cor.~4.16]{Bonf26}}]\label{prop:heart-P}
Let $P$ be $\Pn nk$-like with $k\ge1$ and $n\ge1$, and let $\cH_P$ be as in \cref{prop:tP}.
\begin{enumerate}
\item If $k\ge2$, then $\cH_P=\add P\simeq\modd\kk$.
\item If $k=n=1$, then $\cH_P$ is hereditary and has no non-zero projective objects. If $\thick P$ is algebraic, then $\real\colon\Db(\cH_P)\to\thick P$ is an equivalence.
\item Let $k=1$ and $n\ge2$. Let $t$ span $\Hom(P,P[1])$, and define $Q$ by the triangle
\[P\xrightarrow{a}Q\xrightarrow{b}P\xrightarrow{t}P[1].\]
Then:
\begin{itemize}
\item $Q\in\cH_P$ is a projective cover of $P$ in $\cH_P$, $\End(Q)\cong\kk[x]/(x^2)$, and $\cH_P\simeq\modd\kk[x]/(x^2)$;
\item $\Hom(Q,P)\cong\kk$, and for $l\ne0$ we have $\Hom(Q,P[l])\cong\kk$ if $l=n$ and $\Hom(Q,P[l])=0$ otherwise;
\item $\Hom(Q,Q[1])=0$.
\end{itemize}
\item In cases \textup{(1)} and \textup{(3)}, the comparison maps $\Ext^\bullet_{\cH_P}(P,P)\to\End^\bullet(P)$ are not bijective. In particular, no realisation functor $\Db(\cH_P)\to\thick P$ is an equivalence.
\end{enumerate}
\end{proposition}

\begin{proof}
Recall that $\Ext^1_{\cH_P}(X,Y)=\Hom(X,Y[1])$ for $X,Y\in\cH_P$ (\cref{prop:ext-comparison}).

(1) We have $\Ext^1_{\cH_P}(P,P)=\End^1(P)=0$ since $k\ge2$. Hence every object of $\Filt(P)$ is isomorphic to some $P^{\oplus r}$, and $\Hom(P,-)$ gives $\cH_P\simeq\modd\kk$.

(2) We have $\Ext^1_{\cH_P}(P,P)\cong\kk$. By \cref{prop:ext-comparison}, $\Ext^2_{\cH_P}(P,P)$ embeds into $\End^2(P)$, which is $0$ for $(k,n)=(1,1)$.
\begin{itemize}
\item By d\'evissage (long exact sequences in both variables), $\Ext^2_{\cH_P}$ vanishes identically. Since Yoneda extensions of length $\ge2$ factor through Yoneda products with $\Ext^2$, all $\Ext^{\ge2}_{\cH_P}$ vanish, so $\cH_P$ is hereditary.
\item Consequently $\chi(X,Y):=\dim\Hom(X,Y)-\dim\Ext^1(X,Y)$ is additive in both variables. It is therefore given by $\chi(X,Y)=\ell(X)\ell(Y)\chi(P,P)=0$, where $\ell$ denotes length.
\item If $R$ is projective, then $\dim\End(R)=\chi(R,R)=0$, so $R=0$.
\item The comparison maps $\Ext^l_{\cH_P}(P,P)\to\End^l(P)$ are bijective for $l\le1$ by \cref{prop:ext-comparison}, and both sides vanish for $l\ge2$. By d\'evissage and the five lemma, $\Ext^l_{\cH_P}(X,Y)\to\Hom(X,Y[l])$ is bijective for all $X,Y\in\cH_P$ and all $l$. Thus $\real$ is fully faithful, and it is an equivalence by \cref{thm:realisation}.
\end{itemize}

(3) Since $\cH_P$ is closed under extensions, $Q\in\cH_P$, and $0\to P\to Q\to P\to0$ is a non-split exact sequence in $\cH_P$. Let $W\in\thick P$. Applying $\Hom(-,W)$ to the triangle gives an exact sequence
\[\Hom(P[1],W)\xrightarrow{t^*}\Hom(P,W)\xrightarrow{b^*}\Hom(Q,W)\xrightarrow{a^*}\Hom(P,W)\xrightarrow{t[-1]^*}\Hom(P[-1],W).\]
For $W=P[l]$, the outer maps are, up to sign, multiplication by $t$, namely $\End^{l-1}(P)\to\End^l(P)$ and $\End^l(P)\to\End^{l+1}(P)$. Hence there is a short exact sequence
\[0\to\operatorname{coker}\big(\End^{l-1}(P)\xrightarrow{t}\End^l(P)\big)\to\Hom(Q,P[l])\to\ker\big(\End^l(P)\xrightarrow{t}\End^{l+1}(P)\big)\to0.\]
Since $\deg t=1$ and $t^n\ne0=t^{n+1}$:
\begin{itemize}
\item the cokernel is $\kk$ for $l=0$ and $0$ otherwise;
\item the kernel is $\kk$ (spanned by $t^n$) for $l=n$ and $0$ otherwise.
\end{itemize}
This proves the second item of (3); note that $0\neq n$.

Since $n\ge2$, we have $\Hom(Q,P[1])=0$, i.e.\ $\Ext^1_{\cH_P}(Q,P)=0$. As $P$ is the only simple object, d\'evissage gives $\Ext^1_{\cH_P}(Q,-)=0$, so $Q$ is projective. Since $\Hom(Q,P)\cong\kk$, the top of $Q$ is $P$, and $Q$ is a projective cover of $P$.

Applying $\Hom(Q,-)$ to the triangle gives an exact sequence
\[\Hom(Q,P[-1])\to\Hom(Q,P)\to\End(Q)\to\Hom(Q,P)\to\Hom(Q,P[1]).\]
Its outer terms vanish, so $\dim\End(Q)=2$. The algebra $\End(Q)$ is local because $Q$ has simple top; hence $\End(Q)\cong\kk[x]/(x^2)$. A Hom-finite length category with a projective generator $Q$ is equivalent to $\modd\End(Q)$ via $\Hom(Q,-)$, so $\cH_P\simeq\modd\kk[x]/(x^2)$. Finally, $\Hom(Q,Q[1])=\Ext^1_{\cH_P}(Q,Q)=0$.

(4) In case (1), $\Ext^{nk}_{\cH_P}(P,P)=0$ while $\End^{nk}(P)\ne0$. In case (3), $\Ext^{n+1}_{\kk[x]/(x^2)}(\kk,\kk)\cong\kk$, because the minimal projective resolution of $\kk$ is periodic, while $\End^{n+1}(P)=0$. A realisation functor induces the comparison maps on extension groups, so it cannot be fully faithful.
\end{proof}

\begin{remark}
Parts (1)--(3) overlap with \cite[Lemma~3.11]{Bonf26}, where the indecomposable objects of $\cH_P$ are listed and the spaces $\Hom(X_2,P[i])$ are computed; our $Q$ is Bonfert's $X_2$. They also overlap with \cite[Lemma~4.15]{Bonf26}, which shows that $\cH_P$ has an injective cogenerator unless $n=k=1$. The realisation statement in (2) corresponds to Bonfert's Koszul-dual model \cite[Cor.~4.16(2)]{Bonf26}, which identifies $\thick P$ with the thick subcategory of $\mathcal{D}(\kk[[t]])$ generated by the simple module. What is new in \cref{prop:heart-P} is only the explicit identification $\cH_P\simeq\modd\kk[x]/(x^2)$ in (3) and the reformulation of these facts in terms of the conditions \eqref{eq:star} and \eqref{eq:starstar} used for gluing. We have given the proofs in the form needed below.
\end{remark}

\begin{remark}\label{rem:koszul}
\leavevmode
\begin{enumerate}
\item For $n=1$, parts (1) and (2) are consistent with \cref{thm:HJY}. For $w=1$, $\thick P$ is the derived category of finite-dimensional nilpotent $\kk[x]$-modules; cf.\ \cite[Lemma~3.11(3), Cor.~4.16(2)]{Bonf26}. For $w=2$, part~(1) is exactly the phenomenon of \cref{ex:BB-5.11}.
\item The simple module $\kk$ generates $\Db(\modd\kk[x]/(x^2))$, and $\End^\bullet(\kk)\cong\kk[s]$ with $\deg s=1$. So $\kk$ is a ``$\Pn\infty1$-like'' object (\cref{rem:Bonfert-Koszul} with $n=1$, $k=0$). A $\Pn n1$-like object with $n\ge2$ has the same heart, but its endomorphism algebra is truncated at $t^{n+1}=0$. Meanwhile, the regular module $\kk[x]/(x^2)$ is $\Pn10$-like. Thus the cases $k=0$ and $k=1$ are Koszul dual to each other, which is the mechanism behind \cref{thm:HJY}.
\item The non-vanishing $\Hom(Q,P[n])\ne0$ in (3) and $\End^{k}(P)\ne0$ in (1) are exactly the obstructions to the condition \eqref{eq:starstar} of \cref{thm:A}. The vanishing of $\Hom(Q,P[2])$, i.e.\ $n\ne2$, and of $\End^2(P)$, i.e.\ $k\ne2$, is condition \eqref{eq:star}.
\end{enumerate}
\end{remark}


\subsection{Semi-orthogonal sequences of \texorpdfstring{$\mathbb{P}$}{P}-like objects}\label{sec:P-sequences}


\begin{definition}\label{def:P-sequence}
A \emph{$\mathbb{P}$-sequence} in $\cD$ is a sequence $(P_1,\dots,P_m)$ of objects with the following properties:
\begin{itemize}
\item each $P_s$ is $\Pn{n_s}{k_s}$-like for some $n_s\ge0$ and some $k_s\ge1$ (for $n_s=0$ the value of $k_s$ plays no role);
\item $\Hom^\bullet(P_t,P_s)=0$ for $t>s$.
\end{itemize}
Put $\cD_s:=\thick{P_1,\dots,P_s}$. The sequence is \emph{full} if $\cD_m=\cD$, and \emph{admissible} if $\cD_{s-1}$ is admissible in $\cD_s$ for all $s$.
\end{definition}

If all $n_s=0$, a $\mathbb{P}$-sequence is an exceptional sequence, and admissibility holds automatically under mild finiteness assumptions \cite{Bon89}. For $n_s\ge1$ the subcategory $\thick{P_s}\simeq\operatorname{perf}(A_{n_s,k_s})$ is proper but not of the form $\Db(\kk)$. We do not know a general criterion for admissibility, and therefore impose it as a hypothesis.

\begin{intuition}
A $\mathbb{P}$-sequence is to a $\mathbb{P}$-like object what an exceptional sequence is to an exceptional object: an admissible filtration all of whose strata are generated by a single $\mathbb{P}$-like object. Since every stratum carries, up to shift, a unique non-trivial t-structure (\cref{rem:tP-unique}), there is essentially one way to glue, and the question is what the resulting heart looks like.
\end{intuition}

\begin{lemma}\label{lem:P-strata}
Let $(P_1,\dots,P_m)$ be a full admissible $\mathbb{P}$-sequence in an algebraic triangulated category $\cD$.
\begin{enumerate}
\item The subcategory $\thick{P_s}$ equals ${}^\perp\cD_{s-1}\cap\cD_s$, and the quotient functor induces an equivalence $\thick{P_s}\simeq\cT_s$. In particular $\cT_s\simeq\operatorname{perf}(A_{n_s,k_s})$, and $j_{s!}$ is the inclusion $\thick{P_s}\subset\cD_s$.
\item Put $E_s:=P_s=j_{s!}P_s$ and $F_s:=j_{s*}P_s$. Then $F_s$ is the cone of the counit $i_{s*}i_s^!P_s\to P_s$, i.e.\ the left mutation of $P_s$ through $\cD_{s-1}$. If all $n_s=0$, then $(F_m,\dots,F_1)$ is the left dual exceptional sequence.
\item $\Hom^\bullet(E_s,F_t)\cong\delta_{st}\End^\bullet(P_s)$.
\end{enumerate}
\end{lemma}

\begin{proof}
(1) By admissibility, the quotient functor $j_s^*$ restricts to an equivalence ${}^\perp\cD_{s-1}\cap\cD_s\simeq\cT_s$. The subcategory ${}^\perp\cD_{s-1}\cap\cD_s$ is thick and contains $P_s$, hence contains $\thick{P_s}$. The restriction of $j_s^*$ to $\thick{P_s}$ is fully faithful. Its essential image is thick, because $\thick{P_s}$ is idempotent complete, and it contains $j_s^*P_s$. Since $\cD_s$ is generated by $\cD_{s-1}$ and $P_s$, the object $j_s^*P_s$ generates $\cT_s$. Hence the restriction is an equivalence, and $\thick{P_s}={}^\perp\cD_{s-1}\cap\cD_s$. The model $\operatorname{perf}(A_{n_s,k_s})$ comes from \cref{thm:formality}.

(2) The first claim is the triangle $i_*i^!X\to X\to j_*j^*X$ of \cref{def:recollement} for $X=P_s$. For exceptional sequences this is the description $F_s=j_{s*}(\kk)$ of the left dual recalled in \cref{sec:prelim}.

(3) This is \cref{lem:hom-duality} with the test objects $Q_s=S_s=P_s$.
\end{proof}

We equip each $\cT_s=\thick{P_s}$ with the t-structure $t_{P_s}$ of \cref{prop:tP}. By \cref{rem:tP-unique}, every t-structure on $\cD$ glued from non-trivial t-structures on the strata is obtained from this one by shifting the t-structure on each stratum by some integer $p(s)$ (a ``perversity''). Replacing $P_s$ by $P_s[p(s)]$ does not change the filtration, so it suffices to treat the unshifted case. Let $\cA$ denote the glued heart.

\begin{proposition}[Criterion]\label{prop:P-criterion}
$E_s,F_s\in\cA$ for all $s$ if and only if $\Hom(E_s,E_t[l])=0$ and $\Hom(F_s,F_t[l])=0$ for all $s,t$ and all $l<0$.
\end{proposition}

\begin{proof}
This is \cref{prop:criterion} with the test objects $Q_s=S_s=P_s$; note that $\End^\bullet(P_s)$ is concentrated in non-negative degrees.
\end{proof}

\begin{theorem}\label{thm:Pseq}
Let $(P_1,\dots,P_m)$ be a full admissible $\mathbb{P}$-sequence in an algebraic triangulated category $\cD$, and assume $E_s,F_s\in\cA$ for all $s$.
\begin{enumerate}
\item Suppose that for every $s$ either $n_s=0$ or $k_s\ge3$. Then $\cA$ is a highest weight category with respect to $1<\dots<m$, with standard objects $E_s$ and costandard objects $F_s$.
\item Suppose that for every $s$ one of the following holds: $n_s=0$; or $k_s\ge3$; or $k_s=1$ and $n_s\ge3$. Then $\cA\simeq\modd\Lambda$ with $\Lambda$ standardly stratified with respect to $1<\dots<m$, and:
\begin{itemize}
\item $B_s\cong\kk$ in the first two cases, and $B_s\cong\kk[x]/(x^2)$ in the third case;
\item the standard objects are $\Delta(s)=E_s$ in the first two cases, and $\Delta(s)=j_{s!}Q_s$ in the third case, where $Q_s\in\thick{P_s}$ is the object of \cref{prop:heart-P}(3); in the third case there is an exact sequence $0\to E_s\to\Delta(s)\to E_s\to0$ in $\cA$;
\item the proper standard objects are $\Deltabar(s)\cong E_s$, and the proper costandard objects are $\nablabar(s)\cong F_s$.
\end{itemize}
\item In the situations of \textup{(1)} and \textup{(2)}, $\real\colon\Db(\cA)\to\cD$ is an equivalence if and only if $n_s=0$ for all $s$.
\item If $k_s=n_s=1$ for some $s$, then the simple object $L_s$ of $\cA$ has no projective cover. In particular $\cA$ is not equivalent to $\modd\Lambda$ for any finite-dimensional algebra $\Lambda$.
\end{enumerate}
\end{theorem}

\begin{proof}
\emph{Reduction to \cref{thm:A}.} In the situations of (1) and (2) we are in the projective case of \cref{set:local}, with $S_s=P_s$. Moreover:
\begin{itemize}
\item if $n_s=0$ or $k_s\ge2$, then $\Pi_s=P_s$ and $B_s=\kk$ (\cref{prop:heart-P}(1));
\item if $k_s=1$ and $n_s\ge2$, then $\Pi_s=Q_s$ and $B_s\cong\kk[x]/(x^2)$ (\cref{prop:heart-P}(3)).
\end{itemize}
Condition \eqref{eq:star} holds:
\begin{itemize}
\item for $n_s=0$ trivially;
\item for $k_s\ge3$ because $\End^\bullet(P_s)$ is concentrated in degrees divisible by $k_s$;
\item for $k_s=1$ and $n_s\ge3$ because $\Hom(Q_s,P_s[l])\ne0$ only for $l\in\{0,n_s\}$.
\end{itemize}
Since $\Pi_s\in\Filt(P_s)$, we have $\Delta_s=j_{s!}\Pi_s\in\Filt(E_s)\subset\cA$, and $\nablabar_s=F_s$. Hence \cref{thm:A} applies.

\emph{Proof of (2).} Everything except the statement about $\Deltabar(s)$ in the third case is \cref{thm:A}(1),(2). In the third case, applying $j_{s!}$ to $0\to P_s\to Q_s\to P_s\to0$ gives a triangle $E_s\to\Delta(s)\to E_s\to$ with all vertices in $\cA$, i.e.\ an exact sequence. We have $[\Delta(s):L_s]=2$, since $j_s^*\Delta(s)=Q_s$ has length two, and $[E_s:L_s]=1$. Let $N\subset\Delta(s)$ be a subobject with $[\Delta(s)/N:L_s]=1$ and $N\subset E_s$, where $E_s\subset\Delta(s)$ is the subobject above. Then $E_s/N$ has no composition factor $L_s$, so it lies in $i_{s*}\cA_{s-1}$. But $E_s=j_{s!}P_s={}^p\!j_{s!}P_s$ has no non-zero quotient in $i_{s*}\cA_{s-1}$ (\cref{prop:glued-heart}(2)), so $N=E_s$. Since every quotient of $\Delta(s)$ with a single composition factor $L_s$ factors through $\Delta(s)/E_s$ (because $\rad\End(\Delta(s))$ maps $\Delta(s)$ onto the subobject $E_s$), we conclude $\Deltabar(s)=\Delta(s)/E_s\cong E_s$.

\emph{Proof of (1).} We have $B_s=\kk$ for all $s$, so $\cA$ is quasi-hereditary, i.e.\ highest weight. For quasi-hereditary algebras $\nabla(s)=\nablabar(s)$, so the costandard objects are the $F_s$.

\emph{Proof of (3).} By \cref{thm:A}(3) it suffices to check when \eqref{eq:starstar} holds. If $n_s\ge1$ and $k_s\ge3$, then $\Hom(P_s,P_s[k_s])\ne0$. If $k_s=1$ and $n_s\ge3$, then $\Hom(Q_s,P_s[n_s])\ne0$. If $n_s=0$, the condition holds.

\emph{Proof of (4).} Suppose $R$ is a projective cover of $L_s$ in $\cA$, and let $\bar R$ be its largest quotient lying in the Serre subcategory $\cA_s$; then $\bar R$ is a projective cover of $L_s$ in $\cA_s$.
\begin{itemize}
\item For $M\in\cH_s$ we have $j_{s*}M\in\Filt(F_s)\subset\cA_s$. Hence
\[\Ext^1_{\cH_s}(j_s^*\bar R,M)=\Hom(\bar R,j_{s*}M[1])=\Ext^1_{\cA_s}(\bar R,j_{s*}M)=0,\]
so $j_s^*\bar R$ is projective in $\cH_s$.
\item $j_s^*\bar R\ne0$: otherwise $\bar R\in i_{s*}\cA_{s-1}$, and then $\Hom(\bar R,L_s)=0$ because $L_s=j_{s!*}S_s$ has no non-zero subobject in $i_{s*}\cA_{s-1}$.
\end{itemize}
This contradicts \cref{prop:heart-P}(2).
\end{proof}

\begin{example}[A case with $k=2$]\label{ex:k2}
Let $\Gamma$ be the graded $\kk$-linear category (a dg category with zero differential) with two objects $1,2$ and the following morphism spaces:
\begin{itemize}
\item $\Gamma(1,1)=\kk$ and $\Gamma(2,2)=\kk[t]/(t^2)$ with $\deg t=2$;
\item $\Gamma(1,2)=\kk f$ with $\deg f=0$, and $\Gamma(2,1)=0$;
\item $t\circ f=0$.
\end{itemize}
Let $\cD=\operatorname{perf}(\Gamma)$, and let $P_1,P_2$ be the representable objects, so that $\Hom^\bullet(P_i,P_j)=\Gamma(i,j)$. Then $(P_1,P_2)$ is a full $\mathbb{P}$-sequence with $n_1=0$ and with $P_2$ a $\Pn12$-like object. It is admissible, because $\thick{P_1}\simeq\Db(\kk)$ is generated by an exceptional object and $\cD$ is proper.

We have $i^!P_2=P_1$, hence $F_2=\cone(P_1\xrightarrow{f}P_2)$ and $F_1=E_1=P_1$. Checking the negative Homs:
\begin{itemize}
\item $\Hom^\bullet(E_1,E_2)=\kk f$ is concentrated in degree $0$;
\item $\Hom^\bullet(F_1,F_2)=0$;
\item $\Hom^l(F_2,F_1)\cong\End^{l-1}(P_1)$ is concentrated in degree $1$;
\item $\End^\bullet(F_2)\cong\End^\bullet(P_2)$.
\end{itemize}
By \cref{prop:P-criterion}, $E_s,F_s\in\cA$ for $s=1,2$.

The triangle $P_1\to P_2\to F_2\to P_1[1]$ has all vertices in $\cA$, so it gives an exact sequence $0\to E_1\to E_2\to F_2\to0$. Hence the simple objects are $L_1=E_1$ and $L_2=F_2$. Their extensions are:
\begin{itemize}
\item $\Ext^1_\cA(L_2,L_1)=\Hom(F_2,P_1[1])\cong\kk$;
\item $\Ext^1_\cA(L_1,L_2)=\Hom(P_1,F_2[1])=0$;
\item $\Ext^1_\cA(L_1,L_1)=\Ext^1_\cA(L_2,L_2)=0$.
\end{itemize}
Moreover $\Ext^1_\cA(E_1,-)$ and $\Ext^1_\cA(E_2,-)$ vanish on both simple objects, so $E_1,E_2$ are the indecomposable projective objects. Thus $\cA\simeq\modd\kk\vec A_2$ is highest weight with standard objects $E_1,E_2$ and costandard objects $F_1,F_2$: the conclusion of \cref{thm:Pseq}(1) holds even though $k_2=2$.

This does not contradict the hypotheses of \cref{thm:Pseq}: here $V_1=\Ext^1_\cA(E_1,E_2)=0$, so by \cref{rem:star} condition \eqref{eq:star} is not needed for the top stratum. The realisation functor is not an equivalence, since $\Hom(P_2,P_2[2])\ne0$ while $\Ext^2_{\kk\vec A_2}=0$.
\end{example}

\begin{remark}[The obstruction for $k=2$]\label{rem:obstruction}
Suppose all strata have semisimple hearts, i.e.\ $k_s\ge2$ or $n_s=0$. In the induction step (\cref{prop:step}(2)), the only point where \eqref{eq:star} is needed for an earlier stratum $t<m$ is the case $l=1$ (step (e)). There the kernel $K$ of $\Delta_m\twoheadrightarrow L_m$ lies in $i_*\cA_{m-1}$, and one needs $\Hom(G_s,K[2])=0$. If some $k_t=2$ with $t<m$, this group can be non-zero, and whether $\Ext^1_\cA(G_s,L_m)$ vanishes then depends on the connecting map $\Hom(G_s,L_m[1])\to\Hom(G_s,K[2])$. This map is determined by higher (Massey-type) composition data and not by the graded endomorphism algebras alone; compare the non-formality of the Koszul duals in \cref{rem:Bonfert-Koszul}.
\end{remark}

\begin{question}\label{q:k2}
Let $(P_1,\dots,P_m)$ be a full admissible $\mathbb{P}$-sequence with $E_s,F_s\in\cA$, and suppose some $k_s=2$, or some $(k_s,n_s)=(1,2)$. Is $\cA$ Deligne finite? If so, is it standardly stratified with standard objects $j_{s!}\Pi_s$, as in \cref{thm:Pseq}?
\end{question}

A natural source of test cases is geometry:
\begin{itemize}
\item the structure sheaves of hyperk\"ahler manifolds are $\Pn n2$-objects;
\item by \cite[Thm.~B]{BC25}, as quoted in \cite{Bonf26}, all indecomposable perverse sheaves on $\mathbb{P}^n$ are $\mathbb{P}[2]$-like, and most of them are not Calabi--Yau.
\end{itemize}
It would be interesting to know whether the latter can be arranged into $\mathbb{P}$-sequences satisfying the hypotheses above.

% ===============================================================
\section{Further questions}\label{sec:questions}
% ===============================================================

\begin{question}[Classification of glued t-structures]\label{q:classification}
Let $(P_1,\dots,P_m)$ be a full admissible $\mathbb{P}$-sequence with all $k_s\ge1$. By \cref{thm:Bonfert}, the t-structures on $\cD$ glued along the filtration from non-trivial t-structures on the strata are parametrised by perversity functions $p\colon\{1,\dots,m\}\to\mathbb{Z}$. Which t-structures on $\cD$ are glued in this sense? For which perversities do the objects $E_s[p(s)]$ and $F_s[p(s)]$ lie in the heart, as for the middle perversity in \cite[\S5.1]{BB26}?
\end{question}

\begin{question}[Ringel duality]
By \cite[Thm.~3.2]{BB26}, the Ringel dual of a highest weight category glued along an exceptional sequence is glued along the left dual sequence. Characteristic tilting modules and Ringel duality are available for standardly and properly stratified algebras (see e.g.\ \cite{AHLU00}). Is there a ``dual'' local-exceptional sequence and Hom-dual sequence, built from $(E_\bullet,N_\bullet)$ in \cref{thm:B}, whose gluing produces the Ringel dual?
\end{question}

\begin{question}[The weight-structure side]
In \cref{thm:B}(3), the subcategory $\thick{E_1,\dots,E_m}\simeq\Kb(\proj\cA)$ is filtered with strata $\thick{E_s}\simeq\Kb(\proj B_s)$. These strata carry the co-t-structures with co-hearts $\add E_s$ (\cref{prop:local-exc}). Is the co-t-structure obtained by gluing them \cite{LVY14} the one with co-heart $\add G$? Does it correspond to the glued t-structure under the Koenig--Yang bijection (\cref{thm:KY})? A positive answer would give a description of \cref{thm:B} purely in terms of the local-exceptional sequence and its silting gluing, with the Hom-dual sequence appearing as the Koenig--Yang dual simple-minded collection.
\end{question}

\begin{question}[Without a right dual]
Is there a version of \cref{thm:A,thm:B} assuming only (HD1)--(HD3), i.e.\ only the existence of the functors $j_{s*}$, which is automatic by \cref{lem:sod}? This would require a one-sided gluing of t-structures along semi-orthogonal decompositions (cf.\ \cite{CP10}). It would remove the only hypothesis of \cref{thm:B} that is not automatic.
\end{question}

% ===============================================================
\begin{thebibliography}{BBD82}
% ===============================================================

\bibitem[AHLU00]{AHLU00}
I.~\'Agoston, D.~Happel, E.~Luk\'acs, L.~Unger, \emph{Standardly stratified algebras and tilting}, J.~Algebra \textbf{226} (2000), 144--160.

\bibitem[AI12]{AI12}
T.~Aihara, O.~Iyama, \emph{Silting mutation in triangulated categories}, J.~Lond. Math. Soc. (2) \textbf{85} (2012), 633--668.

\bibitem[AN09]{AN09}
S.~Al-Nofayee, \emph{Simple objects in the heart of a t-structure}, J.~Pure Appl. Algebra \textbf{213} (2009), 54--59.

\bibitem[BB26]{BB26}
A.~Bodzenta, A.~Bondal, \emph{Highest weight categories via pairs of dual exceptional sequences}, arXiv:2601.22004 (2026).

\bibitem[BBD82]{BBD82}
A.~A.~Beilinson, J.~Bernstein, P.~Deligne, \emph{Faisceaux pervers}, Ast\'erisque \textbf{100} (1982).

\bibitem[BC25]{BC25}
L.~Bonfert, A.~Cipriani, \emph{Serre functor and $\mathbb{P}$-objects for perverse sheaves on $\mathbb{P}^n$}, arXiv:2506.06051 (2025).

\bibitem[Bei87]{Bei87}
A.~A.~Beilinson, \emph{On the derived category of perverse sheaves}, in: K-theory, arithmetic and geometry (Moscow, 1984--1986), Lecture Notes in Math. \textbf{1289}, Springer, 1987, 27--41.

\bibitem[BK89]{BK89}
A.~I.~Bondal, M.~M.~Kapranov, \emph{Representable functors, Serre functors, and reconstructions}, Izv. Akad. Nauk SSSR Ser. Mat. \textbf{53} (1989), 1183--1205.

\bibitem[Bon89]{Bon89}
A.~I.~Bondal, \emph{Representations of associative algebras and coherent sheaves}, Izv. Akad. Nauk SSSR Ser. Mat. \textbf{53} (1989), 25--44.

\bibitem[Bon10]{Bon10}
M.~V.~Bondarko, \emph{Weight structures vs.\ t-structures; weight filtrations, spectral sequences, and complexes (for motives and in general)}, J.~K-Theory \textbf{6} (2010), 387--504.

\bibitem[Bonf26]{Bonf26}
L.~Bonfert, \emph{Weight structures and t-structures on the triangulated category generated by a $\mathbb{P}$-object}, arXiv:2610.21859 (2026).

\bibitem[CP10]{CP10}
J.~Collins, A.~Polishchuk, \emph{Gluing stability conditions}, Adv. Theor. Math. Phys. \textbf{14} (2010), 563--607.

\bibitem[CPS88]{CPS88}
E.~Cline, B.~Parshall, L.~Scott, \emph{Finite-dimensional algebras and highest weight categories}, J.~Reine Angew. Math. \textbf{391} (1988), 85--99.

\bibitem[CPS96]{CPS96}
E.~Cline, B.~Parshall, L.~Scott, \emph{Stratifying endomorphism algebras}, Mem. Amer. Math. Soc. \textbf{124} (1996), no.~591.

\bibitem[Dla96]{Dla96}
V.~Dlab, \emph{Quasi-hereditary algebras revisited}, An. \c{S}tiin\c{t}. Univ. Ovidius Constan\c{t}a Ser. Mat. \textbf{4} (1996), 43--54.

\bibitem[DR92]{DR92}
V.~Dlab, C.~M.~Ringel, \emph{The module theoretical approach to quasi-hereditary algebras}, in: Representations of algebras and related topics (Kyoto, 1990), London Math. Soc. Lecture Note Ser. \textbf{168}, Cambridge Univ. Press, 1992, 200--224.

\bibitem[ES03]{ES03}
K.~Erdmann, C.~S\'aenz, \emph{On standardly stratified algebras}, Comm. Algebra \textbf{31} (2003), 3429--3446.

\bibitem[HJY13]{HJY13}
T.~Holm, P.~J\o rgensen, D.~Yang, \emph{Sparseness of t-structures and negative Calabi--Yau dimension in triangulated categories generated by a spherical object}, Bull. Lond. Math. Soc. \textbf{45} (2013), 120--130.

\bibitem[HK19]{HK19}
A.~Hochenegger, A.~Krug, \emph{Formality of $\mathbb{P}$-objects}, Compos. Math. \textbf{155} (2019), no.~5.

\bibitem[HT06]{HT06}
D.~Huybrechts, R.~Thomas, \emph{$\mathbb{P}$-objects and autoequivalences of derived categories}, Math. Res. Lett. \textbf{13} (2006), 87--98.

\bibitem[Kal17]{Kal17}
M.~Kalck, \emph{Derived categories of quasi-hereditary algebras and their derived composition series}, in: Representation theory---current trends and perspectives, EMS Ser. Congr. Rep., Eur. Math. Soc., 2017, 269--308.

\bibitem[Kel94]{Kel94}
B.~Keller, \emph{Deriving DG categories}, Ann. Sci. \'Ec. Norm. Sup\'er. (4) \textbf{27} (1994), 63--102.

\bibitem[KN13]{KN13}
B.~Keller, P.~Nicol\'as, \emph{Weight structures and simple dg modules for positive dg algebras}, Int. Math. Res. Not. IMRN \textbf{2013}, no.~5.

\bibitem[Kru18]{Kru18}
A.~Krug, \emph{Varieties with $\mathbb{P}$-units}, Trans. Amer. Math. Soc. \textbf{370} (2018), no.~11.

\bibitem[KY14]{KY14}
S.~Koenig, D.~Yang, \emph{Silting objects, simple-minded collections, t-structures and co-t-structures for finite-dimensional algebras}, Doc. Math. \textbf{19} (2014), 403--438.

\bibitem[Let24]{Let24}
J.~C.~Letz, \emph{Realization functors in algebraic triangulated categories}, arXiv:2412.07455 (2024).

\bibitem[LVY14]{LVY14}
Q.~Liu, J.~Vit\'oria, D.~Yang, \emph{Gluing silting objects}, Nagoya Math. J. \textbf{216} (2014), 117--151.

\bibitem[Maz04]{Maz04}
V.~Mazorchuk, \emph{On finitistic dimension of stratified algebras}, Algebra Discrete Math. (2004), no.~3, 77--88.

\bibitem[MMS04]{MMS04}
E.~N.~Marcos, O.~Mendoza, C.~S\'aenz, \emph{Stratifying systems via relative simple modules}, J.~Algebra \textbf{280} (2004), 472--487.

\bibitem[Pau08]{Pau08}
D.~Pauksztello, \emph{Compact corigid objects in triangulated categories and co-t-structures}, Cent. Eur. J. Math. \textbf{6} (2008), 25--42.

\bibitem[ST01]{ST01}
P.~Seidel, R.~Thomas, \emph{Braid group actions on derived categories of coherent sheaves}, Duke Math. J. \textbf{108} (2001), 37--108.

\end{thebibliography}

\end{document}
